% DER_XIIc-b: Grothendieck's text with the commentary set in, then the overview, inside a group that applies the chunk's own definitions
\dgrpart{1}{XIIc-b}{Section 6, part 2 of 2, manuscript pages 235--275}{XIIc-b\enspace\textperiodcentered\enspace Chapter XII, section 6, part 2 of 2\enspace\textperiodcentered\enspace manuscript pages 235--275}{Chapter XII \textperiodcentered\ \textup{XIIc-b}}
\dgrchunkbegin
% DER_XIIc-b: the definitions of the chunk's preamble, applied inside its group
\tikzcdset{every diagram/.append style={nodes in empty cells}}
\fancyhf{}
\fancyhead[L]{\footnotesize\itshape Les D\'erivateurs -- XIIc-b}
\fancyhead[R]{\footnotesize Manuscript 228--275}
\fancyfoot[C]{\footnotesize\thepage}
\providecommand\headrulewidth{}\renewcommand{\headrulewidth}{0.3pt}
\let\printedpage\relax\newcommand{\printedpage}[1]{\marginpar{\footnotesize p.~#1}}
\let\msp\relax\newcommand{\msp}[1]{\marginpar{\footnotesize [#1]}}
\let\cX\relax\newcommand{\cX}{\mathscr{X}}
\let\Ob\relax\newcommand{\Ob}{\operatorname{Ob}}
\let\cercle\relax\newcommand{\cercle}[1]{\tikz[baseline=(c.base)]\node[draw,circle,inner sep=1pt](c){#1};}
\let\Fibu\relax\newcommand{\Fibu}{\underline{\mathrm{Fib}}}
\let\Homu\relax\newcommand{\Homu}{\underline{\mathrm{Hom}}}
\let\Ho\relax\newcommand{\Ho}{\mathrm{Ho}}
\let\ho\relax\newcommand{\ho}{\mathrm{ho}}
\let\lc\relax\newcommand{\lc}{\mathrm{lc}}
\let\uFib\relax\newcommand{\uFib}{\underline{\mathrm{Fib}}}
\let\Fib\relax\newcommand{\Fib}{\mathrm{Fib}}
\let\Cof\relax\newcommand{\Cof}{\mathrm{Cof}}
\let\Asph\relax\newcommand{\Asph}{\mathrm{Asph}}

% ---- from Der-ChXIIc-30-03-03_p036-040.tex
\printedpage{150}
% head: A. Grothendieck : Dérivateurs, Ch. XII, Édition des Universités de Montpellier II et Paris VII
% foot: Version 30 mars 2003

Reste à prouver c). Je vais le déduire d'une

\textbf{Proposition 17.} \textit{Soit un carré commutatif}
\[
Q \quad = \quad
\begin{tikzcd}[ampersand replacement=\&]
X_0 \arrow[r, "i_1"] \arrow[d, "i_2"'] \& X_1 \arrow[d, "f_1"] \\
X_2 \arrow[r, "f_2"] \& Y
\end{tikzcd}
\]
\textit{dans} Cat\textit{, et $\cX_Q$ la catégorie fibrée sur}
\[
\tilde{\Phi} \quad (\simeq \mathbf{\Delta}^1 \times \mathbf{\Delta}^1) \quad = \quad
\begin{tikzcd}[ampersand replacement=\&]
0 \& 1 \arrow[l] \\
2 \arrow[u] \& 3 \arrow[l] \arrow[u]
\end{tikzcd}
\]
\textit{qui l'intègre. Soit $\cX$ l'image inverse du fermé}
\[
\begin{tikzcd}[ampersand replacement=\&]
0 \& 1 \arrow[l] \\
2 \arrow[u] \&
\end{tikzcd}
\]
\msp{page 235}
\textit{dans $\cX_Q$, qui intègre le diagramme}
\[
\begin{tikzcd}[ampersand replacement=\&, row sep=small]
 \& X_1 \\
X_0 \arrow[ru] \arrow[rd] \& \\
 \& X_2
\end{tikzcd}
\]
\textit{. Pour que $Q$ soit $W$-cocartésien, il faut et il suffit que l'inclusion}
\[
i : \cX \longrightarrow \cX_Q
\]
\textit{soit dans $W$.}

Comme $\cX_Q$ est lisse sur $\tilde{\Phi}$, et 3 est objet initial de $\tilde{\Phi}$, il s'ensuit que l'inclusion de la fibre de $\cX_Q$ en le sommet 3 de $\tilde{\Phi}$ est dans $W$,
\[
Y \xrightarrow{i_Y} \cX_Q \qquad\qquad \in W.
\]

Considérons alors le diagramme
\[
\begin{tikzcd}[ampersand replacement=\&, column sep=small]
\cX \arrow[rr, "\varphi_Q"] \arrow[rd, "i"', ""{name=direct,pos=.55,above right}] \& \& Y \arrow[ld, "i_Y\in W", ""{name=via,pos=.30,below left}] \\
 \& \cX_Q. \& \arrow[from=via,to=direct,dotted,shorten <=1pt,shorten >=1pt,"\alpha"']
\end{tikzcd}
\]

\printedpage{151}
% head: A. Grothendieck : Dérivateurs, Ch. XII, Édition des Universités de Montpellier II et Paris VII
% foot: Version 30 mars 2003

Ce diagramme n'est pas commutatif, mais il y a une flèche de commutativité
\[
\alpha : i_Y\varphi_Q \longrightarrow i.
\]

\dgrblock{dgrb580}{The integral criterion with its natural transformation}
\begin{dcomment}{exposition}{p235:Ce diagramme n'est pas commutatif}{The integral criterion with its natural transformation}
Proposition~17 turns the comparison with a homotopy pushout into a test on one inclusion. This is what makes the pasting proof manageable: several comparison maps can then be placed in the same category and related by literal inclusions. The natural transformation in the triangle is essential to that change of viewpoint.

Here is an explicit model, fixing the variance. Regard the square as a contravariant diagram $F:\widetilde\Phi^{\mathrm{op}}\to\mathrm{Cat}$, with $F(0)=X_0$, $F(1)=X_1$, $F(2)=X_2$, and $F(3)=Y$. The fibred Grothendieck construction $E=\int F$ has objects $(b,x)$ for $x\in F(b)$. A morphism $(b,x)\to(c,y)$ consists of an arrow $u:b\to c$ in the base together with an arrow $x\to F(u)(y)$ in $F(b)$. Composition uses the contravariance of $F$.

The vertex $3$ is initial in $\widetilde\Phi$. Its unique arrows $3\to b$ therefore define a functor
\[
 r:E\longrightarrow Y,
 \qquad r(b,x)=F(3\to b)(x).
\]
For $y\in Y$ there is a natural bijection
\[
 \operatorname{Hom}_E\bigl((3,y),(b,x)\bigr)
   =\operatorname{Hom}_Y\bigl(y,r(b,x)\bigr).
\]
Thus the fibre inclusion $i_Y:Y\to E$ is left adjoint to $r$, with $ri_Y=1_Y$. The counit $i_Yr\Rightarrow1_E$ is the arrow over $3\to b$ represented by the identity of $F(3\to b)x$. This displays the homotopy equivalence of $Y$ with $E$ directly.

Now restrict to $\mathcal X=E|_{\{0,1,2\}}$. The restricted functor $r|_{\mathcal X}$ is exactly the source's comparison $\varphi_Q:\mathcal X\to Y$. Restricting the counit gives
\[
 \alpha:i_Y\varphi_Q\Longrightarrow i,
 \qquad i:\mathcal X\hookrightarrow E.
\]
This is an arrow between functors with source $\mathcal X$ and target $E$. In the edition it is drawn inside the triangle. The transcription had instead attached it to the vertices $Y$ and $\mathcal X$; the correction register restores the interior arrow.

To finish the criterion, a natural transformation $g\Rightarrow h$ determines a functor $\mathcal X\times[1]\to E$. The projection $\mathcal X\times[1]\to\mathcal X$ and its two sections belong to the localizer. Consequently $g$ belongs to $W$ exactly when $h$ does. Applying this to $\alpha$, and then using $i_Y\in W$, gives
\[
 i\in W
 \ \Longleftrightarrow\ i_Y\varphi_Q\in W
 \ \Longleftrightarrow\ \varphi_Q\in W.
\]
The triangle need not commute strictly. Its specified transformation is precisely enough because the localizer already identifies naturally homotopic functors.
\end{dcomment}


On sait que ceci implique que
\[
i \in W \qquad\qquad \Longleftrightarrow \qquad\qquad (i_Y\varphi_Q) \in W
\]
(cf. lemme p. 2), et comme $i_Y \in W$,
\[
(i_Y\varphi_Q) \in W \qquad\qquad \Longleftrightarrow \qquad\qquad \varphi_Q \in W,
\]
d'où
\[
i \in W \qquad\qquad \Longleftrightarrow \qquad\qquad \varphi_Q \in W,
\]

\hfill q.e.d.

Nous pouvons maintenant prouver

\msp{page 236}

\textbf{Proposition 18} (\footnote[127]{devrait être vrai dans toute catégorie de modèles ayant assez de cofibrations.}). (Moyennant W(7 bis).) \textit{Considérons un diagramme commutatif}
\[
\begin{tikzcd}[ampersand replacement=\&]
X_0 \arrow[r] \arrow[d] \arrow[rd, phantom, "1"] \& X_1 \arrow[r] \arrow[d] \arrow[rd, phantom, "2"] \& X_2 \arrow[d] \\
X'_0 \arrow[r] \& X'_1 \arrow[r] \& X'_2
\end{tikzcd}
\tag*{$(\ast)$}
\]
\textit{d'où un carré composé}
\[
\begin{tikzcd}[ampersand replacement=\&]
X_0 \arrow[r] \arrow[d] \arrow[rd, phantom, "3"] \& X_2 \arrow[d] \\
X'_0 \arrow[r] \& X'_2.
\end{tikzcd}
\]
\textit{Supposons que le carré} 1 \textit{soit $W$-cocartésien. Alors le carré} 2 \textit{est $W$-cocartésien si et seulement si le carré} 3 \textit{l'est.}

\textsc{Démonstration.}

\texttt{[Maintenant, }$X_i$\texttt{ et }$X'_i$\texttt{ échangent leurs rôles.]}

\printedpage{152}
% head: A. Grothendieck : Dérivateurs, Ch. XII, Édition des Universités de Montpellier II et Paris VII
% foot: Version 30 mars 2003

Soit
\[
I \quad = \quad
\left(
\begin{tikzcd}[ampersand replacement=\&]
0 \arrow[d] \& 1 \arrow[l] \arrow[d] \& 2 \arrow[l] \arrow[d] \\
0' \& 1' \arrow[l] \& 2' \arrow[l]
\end{tikzcd}
\right)
\]
et soit $\cX$ l'intégrale du diagramme $(\ast)$, catégorie fibrée sur $\cX$. On désigne par $\cX_{010'}$ la sous-catégorie pleine de $\cX$ image inverse de la sous-catégorie pleine % [sic] fibrée sur $\cX$ -> fibrée sur $I$, a category fibred over itself; the base is the category $I$ just defined
\[
\underbrace{I_{010'}}_{=\{0,1,0'\}} \ \subseteq I,
\]
et de même pour toute autre partie $A$ de $\Ob I = \{0,1,2,0',1',2'\}$, donnant une sous-catégorie pleine $\cX_A$ de $\cX$. On pose
\[
U \quad = \quad
\begin{tikzcd}[ampersand replacement=\&]
0 \arrow[d] \& 1 \arrow[l] \& 2 \arrow[l] \\
0' \& \&
\end{tikzcd}
, \qquad
V \quad = \quad
\begin{tikzcd}[ampersand replacement=\&]
0 \arrow[d] \& 1 \arrow[l] \arrow[d] \& 2 \arrow[l] \\
0' \& 1' \arrow[l] \&
\end{tikzcd}
,
\]
d'où
\[
\cX_U = \cX_{0120'}, \qquad\qquad \cX_V = \cX_{0120'1'}.
\]
Les flèches entre les $\cX_A$ sont les inclusions

\msp{page 237}

déduites d'inclusions $A \hookrightarrow A'$ dans $\mathfrak{P}(\Ob I)$. Ainsi, le caractère $W$-cocartésien des carrés 1, 2, 3 s'exprime respectivement par la condition que les inclusions
\[
\begin{aligned}
\cX_{010'} \quad &\xrightarrow{\varphi_1} \quad \cX_{010'1'} \\
\cX_{121'} \quad &\xrightarrow{\varphi_2} \quad \cX_{121'2'} \\
\cX_{020'} \quad &\xrightarrow{\varphi_1} \quad \cX_{020'2'} % [sic] $\varphi_1$ -> $\varphi_3$, this inclusion is called $\varphi_3$ in the diagram that follows
\end{aligned}
\]
soient dans $W$. Considérons alors le diagramme de flèches d'inclusion
\printedpage{153}
% head: A. Grothendieck : Dérivateurs, Ch. XII, Édition des Universités de Montpellier II et Paris VII
% foot: Version 30 mars 2003
\[
\begin{tikzcd}[ampersand replacement=\&]
\cX_{020'} \arrow[r, "\varphi_3"] \arrow[d, "\alpha"'] \& \cX_{020'2'} \arrow[d, "\beta"] \\
\cX_U \arrow[r, "\psi_3"] \arrow[d, "\psi_1"'] \& \cX \arrow[d, "\mathrm{id}"] \\
\cX_V \arrow[r, "\psi_2"] \& \cX \\
\cX_{121'} \arrow[r, "\varphi_2"] \arrow[u, "\alpha'"] \& \cX_{121'2'}. \arrow[u, "\beta'"']
\end{tikzcd}
\]

Je dis que dans ce diagramme on a
\[
\alpha, \beta, \alpha', \beta' \in W
\tag*{$(\ast)$}
\]
et
\[
\varphi_1 \in W \ \ \text{(i.e. carré 1 $W$-cocartésien)} \qquad\qquad \Longleftrightarrow \qquad\qquad \psi_1 \in W.
\tag*{$(\ast\ast)$}
\]

Les relations $(\ast)$ impliquent déja que % [sic] déja -> déjà, missing grave accent
\[
\begin{aligned}
\varphi_2 \in W \quad &\Longleftrightarrow \quad \psi_2 \in W \\
\varphi_3 \in W \quad &\Longleftrightarrow \quad \psi_3 \in W,
\end{aligned}
\]
et sous l'hypothèse $W$-cocartésienne $(\ast\ast)$, impliquant $\psi_1 \in W$,

\msp{page 238}

on trouve
\[
\psi_2 \in W \qquad\qquad \Longleftrightarrow \qquad\qquad \psi_3 \in W,
\]
d'où
\[
\varphi_2 \in W \qquad\qquad \Longleftrightarrow \qquad\qquad \varphi_3 \in W,
\]
ce qu'il fallait prouver.

Il reste à établir $(\ast)$ et $(\ast\ast)$. Les relations $(\ast)$ proviennent de la remarque que si $i : A \longrightarrow A'$ est une inclusion de sous-catégories pleines de $I$, telle que $i$ soit $W$-\textit{asphérique} (i.e. $A/j$ $W$-asphérique pour $j \in \Ob A'\backslash\Ob A$), alors $\cX_A \longrightarrow \cX_{A'}$ l'est aussi (changement de base lisse), donc est \texttt{[dans]} $W$. On doit vérifier la $W$-asphéricité pour les quatre inclusions
\[
\begin{array}{rcll}
\{0,2,0'\} & \longrightarrow & U = \{0,1,2,0'\} & \qquad\qquad \mathbf{\Delta}^0 \\[1ex]
\{0,2,0',2'\} & \longrightarrow & I & \qquad\qquad \mathbf{\Delta}^0/\mathbf{\Delta}^1 \\[1ex]
\{1,2,1'\} & \longrightarrow & V = \{0,1,2,0',1'\} & \qquad\qquad \mathbf{\Delta}^1,\ \Phi \\[1ex]
\{1,2,1',2'\} & \longrightarrow & I & \qquad\qquad \mathbf{\Delta}^1,\ \mathbf{\Delta}^1 \times \mathbf{\Delta}^1.
\end{array}
\]

\dgrblock{dgrb581}{The seven slices behind the pasting theorem}
\begin{dcomment}{examples}{p238:On doit vérifier la $W$-asphéricité}{The seven slices behind the pasting theorem}
The six-object poset is the bookkeeping device that puts the left square, right square and outer rectangle in one ambient diagram. Grothendieck reduces the substantial part of pasting to seven very small slice categories. Computing them explicitly both explains the proof and corrects one harmless shape label in its margin.

Use the arrows of $I$ exactly as printed: the top row is $2\to1\to0$, the bottom row is $2'\to1'\to0'$, and the vertical arrows are $k\to k'$. A slice $A/j$ has as objects the elements of $A$ admitting an arrow to $j$. Since $I$ is a poset, its morphisms are just the induced order. For the inclusions underlying $\alpha,\beta,\alpha',\beta'$ we obtain
\[
\begin{array}{c|c|l}
 \text{inclusion}&j&A/j\\ \hline
 \alpha&1&\{2\}\\
 \beta&1&\{2\}\\
 \beta&1'&(2\longrightarrow2')\\
 \alpha'&0&(2\longrightarrow1)\\
 \alpha'&0'&(2\longrightarrow1\longrightarrow1')\\
 \beta'&0&(2\longrightarrow1)\\
 \beta'&0'&\{1,2,1',2'\}\cong[1]\times[1].
\end{array}
\]
Each displayed category has initial object $2$. For an object $j$ already in $A$, the slice $A/j$ has terminal object $1_j$, so no separate check is necessary. In the fifth line the composite $2\to1'$ exists: the inclusion is full, and the defining rectangle is commutative. The slice is consequently $[2]=\Delta^2$. The edition prints $\Phi$ in this place. The asphericity argument uses the initial object and is unaffected by this label.

We can also make the passage from base slices to total-category slices explicit. Let $F:I^{\mathrm{op}}\to\mathrm{Cat}$ be the integrated diagram and let $e=(j,y)$ be an object of its total category. There is a projection
\[
 (\mathcal X_A)/e\longrightarrow A/j.
\]
An object above $u:a\to j$ is an arrow $x\to F(u)y$ in $F(a)$. The choice $x=F(u)y$, with the identity arrow, defines a section. For an object $z$ of the total slice, every arrow from its image in $A/j$ to $u$ has a unique lift from $z$ to this chosen object over $u$. Hence the projection is left adjoint to the section. The two categories are homotopy equivalent. Asphericity of $A/j$ therefore proves asphericity of $(\mathcal X_A)/e$; the localizer's Theorem~A condition gives $\mathcal X_A\to\mathcal X_{A'}\in W$. This is the ``smooth change of base'' step in concrete form.

The last ingredient is gluing. The two subsets
\[
 U=\{0,1,2,0'\},\qquad
 U'=\{0,1,0',1'\}
\]
are cosieves in $V=\{0,1,2,0',1'\}$: following an arrow from either subset never leaves that subset. Their total categories are therefore closed subcategories covering $\mathcal X_V$, with intersection $\mathcal X_{010'}$. The forward part of the closed-cover axiom sends $\varphi_1\in W$ to $\psi_1\in W$. As $\psi_3=\psi_2\psi_1$, two-out-of-three then says that $\psi_2\in W$ exactly when $\psi_3\in W$. The four slice equivalences transport this assertion back to $\varphi_2$ and $\varphi_3$.

The proof thus needs only the forward part of W(7 bis). The source writes the stronger equivalence $\varphi_1\in W\Longleftrightarrow\psi_1\in W$, which its full hypothesis permits, but its reverse direction is not used. Cancellation of the right square against the outer rectangle is produced by a known weak equivalence $\psi_1$ and two-out-of-three. It should be kept separate from the stronger conservative gluing condition examined in the appendix.
\end{dcomment}


\printedpage{154}
% head: A. Grothendieck : Dérivateurs, Ch. XII, Édition des Universités de Montpellier II et Paris VII
% foot: Version 30 mars 2003

(NB J'ai marqué en marge les catégories qu'on obtient comme $A/j$, à isomorphisme près, dans les quatre cas envisagés - elles ont toutes un objet initial, donc sont $W$-asphériques.)

Reste à prouver $(\ast\ast)$. Soit
\[
U' = \{0,1,0',1'\}, \qquad\qquad \text{donc } V = U \cup U',\ \ U \cap U' = \{0,1,0'\},
\]
où $U$, $U'$ sont des \textit{fermés} de $V$ qui recouvrent $V$. Donc on aura de même
\[
\cX_V = \cX_U \cup \cX_{U'}, \qquad\qquad \cX_U \cap \cX_{U'} = \cX_{010'},
\]

\msp{page 239}

\[
\begin{tikzcd}[ampersand replacement=\&]
\cX_{U\cap U'} \arrow[r] \arrow[d, "\varphi_1"] \& \cX_U \arrow[d, "\psi_1"] \\
\cX_{U'} \arrow[r] \& \cX_V.
\end{tikzcd}
\]

Donc par W(7 bis) (cas dual des immersions fermées), on a
\[
\varphi_1 \in W \qquad\qquad \Longleftrightarrow \qquad\qquad \psi_1 \in W,
\]

\hfill q.e.d.

Notons pour mémoire le corollaire qui nous avait motivé.

\textbf{Corollaire.} \textit{Supposons} W(7 bis) (comme dans toute cette section). \textit{Soit un carré cocartésien dans} Cat
\[
\begin{tikzcd}[ampersand replacement=\&]
X_0 \arrow[r, "i_1"] \arrow[d, "i_2"'] \& X_1 \arrow[d] \\
X_2 \arrow[r] \& X_3,
\end{tikzcd}
\]
\textit{avec $i_1$ (ou $i_2$) une $W$-cofibration (p.ex. une immersion de Dwyer). Alors ce carré est $W$-cocartésien.}

\dgrblock{dgrb582}{From categorical gluing to coherent cocartesian squares}
\begin{dcomment}{later}{p239:Alors ce carré est $W$-cocartésien.}{From categorical gluing to coherent cocartesian squares}
The corollary finishes the comparison that motivated this detour. Factor the second leg of a strict pushout as a Dwyer map followed by a weak equivalence. The first resulting square is homotopy cocartesian by the Dwyer case; the second is so because its opposite sides are weak equivalences, using the first leg's $W$-cofibration property. Proposition~18 identifies their composite as homotopy cocartesian. The categorical integral and the replacement construction therefore give the same answer on these pushouts.

In the modern derivator language, put $\square=[1]\times[1]$ and let $\ulcorner\subset\square$ omit $(1,1)$. A coherent square $Q\in\mathbb D(\square)$ is cocartesian when the counit
\[
 \iota_!\iota^*Q\longrightarrow Q,
 \qquad \iota:\ulcorner\hookrightarrow\square,
\]
is an isomorphism. The composition and cancellation rule for such squares is \citet[Proposition~3.13]{Groth2013}. It requires neither stability nor a zero object. This identifies the modern form of the assertion in footnote~127: the coherence belongs to an object of $\mathbb D([2]\times[1])$, and the left square controls equivalence of the two remaining cocartesianness tests.

For Dwyer maps the later result is stronger still. \citet[Theorems~1.5--1.6, arXiv version~3]{HackneyOzornovaRiehlRovelli2024} show that if $D=B\amalg_A C$ and $A\to B$ is Dwyer, then
\[
 NB\amalg_{NA}NC\longrightarrow ND
\]
is a weak categorical equivalence in the Joyal model structure. Thus the ordinary pushout represents the pushout of the corresponding $(\infty,1)$-categories, retaining categorical information before passage to the homotopy types of their classifying spaces. The hypotheses concern the attaching map, so the result applies directly to the finite cylinders above.

The source's final sentence on cofibres should be read at this finite level. It has exhibited functorial replacements, proved the relative invariance needed to compare them, and established the basic gluing calculus. Those constructions supply the material for the abstract homotopy-square theory in the next section. A claim that all homotopy Kan extensions exist for an arbitrary relative category would require further hypotheses; no such claim is established by this closing factorization alone.
\end{dcomment}


Cette digression (p. 232 ff.) n'est pas nécessaire d'ailleurs pour formuler dès maintenant la conclusion de nos cogitations sur les cofibres $W$-homotopiques : \textit{Il y a une telle théorie dans} Cat, pourvu seulement que $W$ satisfasse à W(7 bis). En effet, tout résulte du fait que toute flèche $f : X \longrightarrow Y$ dans Cat se factorise en
\[
X \xrightarrow{i_f} \tilde{Y} \xrightarrow{p_f} Y,
\]
avec $i_f$ une $W$-cofibration (et même une imm-

% ---- from Der-ChXIIc-30-03-03_p041-045.tex
\printedpage{155}
% head: A. Grothendieck : Dérivateurs, Ch. XII, Édition des Universités de Montpellier II et Paris VII
% foot: Version 30 mars 2003

\msp{page 240}
ersion ouverte de Dwyer), et $p_f \in W$. % beginning of a word continued from the preceding page (outside this range), as printed

\section*{7\quad Retour sur l'axiomatique de la fibre $W$-homotopique.}

Je reviens sur la réflexion du n$^\circ$ 4. Finalement, pour
\[
(\mathcal{M}, W \subseteq \mathrm{Fl}(\mathcal{M}))
\]
donnés avec $W$ faiblement saturé, il apparaît qu'une excellente théorie des fibres $W$-homotopiques, satisfaisant aux conditions \cercle{a}\cercle{b}\cercle{c} du n$^\circ$ 4 (pages 215-218), existe dès que l'axiome d'apparence anodine suivant est satisfait (en plus de l'existence de l'objet final et des produits fibrés (\footnote[128]{et encore! il me semble qu'il suffit de supposer que les $X \longrightarrow Y$ qui sont $W$-fibrants sont `quarrables'.})).

\textbf{Factorisation gauche :} (\footnote[129]{Sous ces conditions, $\mathcal{M}$ muni de $(W, F_W)$ est une catégorie à cofibrations de K.S. \textsc{Brown-Anderson}, et `tout est O.K.', mais il faudrait vérifier a) b) c) ci-dessus (p. 215, 218), je ne suis pas sûr de l'avoir fait.},\footnote[130]{\textbf{NB} L'axiome de factorisation gauche n'est suffisant pour l'existence d'une bonne théorie des fibres homotopiques (semble-t-il) que dans le cas où tout $X$ dans $\mathcal{M}$ est $W$-fibrant, i.e. est dans $\Fibu_W e$.}). Toute flèche $f : X \longrightarrow Y$ de $\mathcal{M}$ se factorise en
\[
X \xrightarrow{i_f} \bar{X} \xrightarrow{p_f} Y,
\]
avec
\[
i_f \in W, \qquad p_f \in \mathrm{Fib}_W.
\]

L'axiome dual donne lieu à une bonne théorie des cofibres $W$-homotopiques.

\textbf{Factorisation droite :} Toute flèche $f : X \longrightarrow Y$ dans $\mathcal{M}$ se factorise en
\[
X \xrightarrow{j_f} \tilde{Y} \xrightarrow{q_f} Y,
\]
avec
\[
q_f \in W, \qquad j_f \in \mathrm{Cofib}_W.
\]

\msp{page 241}
Je voudrais à présent développer quelque peu le formalisme des fibres $W$-homotopiques, sous l'hypothèse (disons) que $W$ satisfait l'axiome de factorisation gauche.

\dgrblock{dgrb583}{The factorization axiom and Brown's hypothesis}
\begin{dcomment}{context}{p240}{The factorization axiom and Brown's hypothesis}
Grothendieck now wants to detach the theory of homotopy fibres from the particular geometry of categories. The interval and mapping-cylinder arguments of the preceding pages have supplied a mechanism; he asks which part of that mechanism belongs to an arbitrary category with weak equivalences. His phrase ``catégorie de modèles'' here means a category equipped with a class $W$ and the properties under discussion. It does not silently import all the axioms of a Quillen model category.

The relevant meaning of $W$-fibration is likewise intrinsic to $W$. Precisely, for $p:X\to Y$, the pullback functor $p^*:\mathcal M/Y\to\mathcal M/X$ must preserve the arrows whose underlying maps belong to $W$. This also implies the same assertion after another base change. This is the property used below, together with the stability of these maps under composition and pullback. In particular, one must not replace $\mathrm{Fib}_W$ by the fibrations of an arbitrary model structure without checking the corresponding properness statement. Here ``weakly saturated'' means that $W$ contains the isomorphisms and satisfies two-out-of-three; it does not mean that every arrow inverted by localization already belongs to $W$. A factorization $f=pi$ with $i\in W$ and $p\in\mathrm{Fib}_W$ is sufficient for the replacement operations that are actually written down; it is not, by its form alone, a pair of lifting axioms.

The two footnotes beneath the axiom record a real issue in its scope. Brown's category of fibrant objects requires every map to the terminal object to be a fibration. Brown also requires pullbacks of fibrations and of trivial fibrations to exist and to remain of the same kind, and a path-object factorization of every diagonal. His factorization lemma produces a factorization of every map into a weak equivalence followed by a fibration; conversely, applying such a factorization to the diagonal supplies the path-object axiom \citep[Part~I, \S1, pp.~420--422]{Brown1973}. Thus the all-object-fibrancy qualification in footnote~130 gives the clean connection with Brown. The word ``cofibrations'' in footnote~129 is a slip for the fibration version of this theory.

Here is why a weakened existence condition on pullbacks is reasonable. If $S$ is fixed and $A\to S$, $B\to S$ are $W$-fibrations, their product over $S$ is formed by pulling back one of those fibrations. The projection to the other factor is again a fibration, and the composite to $S$ is a fibration. The same observation constructs a product with any finite number of factors; the empty product is $S\to S$. A factorization of a diagonal then provides the path objects required in this relative category. Arbitrary pullbacks in $\mathcal M$ are consequently stronger than what these constructions use. Footnote~128 is a proposed economy of axioms, rather than an obstruction to the argument.

The distinction between Brown's category and a category merely admitting fibrant replacement matters for the footnote's present status. In the language of left derivable categories, the written factorization and pullback hypotheses already suffice for the finite diagram theorem. We verify this below, including the required stability of trivial fibrations; no global-fibrancy hypothesis has to be added. Functorial factorization gives a shorter separate proof of the two replacement lemmas, but is also unnecessary in the derivable-category argument. Thus the source's cautious necessity suggestion is not part of the modern theorem.
\end{dcomment}


Considérons un diagramme cartésien
\[
\begin{tikzcd}[ampersand replacement=\&]
X_1 \arrow[d, "p_1"'] \& Y \arrow[l, "q_2"'] \arrow[d, "q_1"] \\
S \& X_2 \arrow[l, "p_2"]
\end{tikzcd}
\]
\printedpage{156}
% head: A. Grothendieck : Dérivateurs, Ch. XII, Édition des Universités de Montpellier II et Paris VII
% foot: Version 30 mars 2003
dans $\mathcal{M}$. Si $p_1$ est dans $\mathrm{Fib}_W$, alors $Y$ sur $X_2$ peut s'interpréter comme la fibre $W$-homotopique de $X_1/S$ en $X_2/S$. Plus précisément, il en est ainsi de la classe $\ho^{\lc}_{W,X_2}(Y)$,
\[
\ho^{\lc}_{W,X_2}(Y) \simeq p_2^*(\ho^{\lc}_{W,S}(X_1/S)). \tag*{$(*)$}
\]
Mais en tant qu'objet sur $X_1$, $Y$ peut s'interpréter également comme fibre $W$-homotopique de $X_2/S$ en $X_1/S$ :
\[
\ho^{\lc}_{W,X_1}(Y/X_1) \simeq p_1^*(\ho^{\lc}_{W,S}(X_2/S)). \tag*{$(**)$}
\]
À dire vrai, il faudrait le prouver. Pour ceci, considérons une factorisation
\[
X_2 \xrightarrow{i_2} \bar{X}_2 \xrightarrow{\bar{p}_2} S \qquad i_2 \in W\ p_2 \in \mathrm{Fib}_W, % [sic] p_2 \in \mathrm{Fib}_W -> \bar{p}_2 \in \mathrm{Fib}_W, the factorisation just written is through \bar{p}_2
\]
de sorte qu'on trouve un diagramme de carrés cartésiens
\[
\begin{tikzcd}[ampersand replacement=\&]
X_1 \arrow[d] \& \bar{Y} \arrow[l, "\bar{q}_2"'] \arrow[d] \& Y \arrow[l, "j_2"', "W"] \arrow[d] \\
S \& \bar{X}_2 \arrow[l, "\bar{q}_2"'] \& X_2. \arrow[l, "i_2"', "W"] % [sic] \bar{q}_2 (bottom left arrow) -> \bar{p}_2, the arrow \bar{X}_2 \to S is \bar{p}_2 of the factorisation above
\end{tikzcd}
\]

\msp{page 242}
Par définition de $p_1^*$, on a
\[
p_1^*(\ho^{\lc}_{W,S}(X_2/S)) \simeq \ho^{\lc}_{W,X_1}(\bar{Y}/X_1),
\]
d'autre part, comme $X_1 \longrightarrow S$ est dans $\mathrm{Fib}_W$, et $i_2 \in W$, on a $j_2 \in W$, donc
\[
\ho^{\lc}_{W,X_1}(Y/X_1) \xrightarrow{\sim} \ho^{\lc}_{W,X_1}(\bar{Y}/X_1),
\]
d'où la formule $(**)$.

Ainsi, l'objet $Y$ a deux interprétations homotopiques différentes, suivant qu'on le regarde comme étant sur $X_1$, ou sur $X_2$. On aimerait alors une interprétation qui englobe les deux à la fois. On va considérer le diagramme
\[
\begin{tikzcd}[ampersand replacement=\&]
X_1 \arrow[d] \& \\
S \& X_2 \arrow[l]
\end{tikzcd}
\]
de type $\Psi^\circ$, ou
\[
\Psi \quad = \quad
\begin{tikzcd}[ampersand replacement=\&]
1 \& \\
0 \arrow[u] \arrow[r] \& 2
\end{tikzcd}
\qquad \text{ou} \qquad
\begin{tikzcd}[ampersand replacement=\&, row sep=small, column sep=small]
 \& 1 \\
0 \arrow[ur] \arrow[dr] \& \\
 \& 2
\end{tikzcd}
\quad ,
\]
% the last "2" of the first diagram is printed overlapping the page number 156
\printedpage{157}
% head: A. Grothendieck : Dérivateurs, Ch. XII, Édition des Universités de Montpellier II et Paris VII
% foot: Version 30 mars 2003
dont $Y$ se déduit par produit fibré, en complétant ce diagramme en un diagramme de type $\tilde{\Psi}^\circ$, où
\[
\tilde{\Psi} =
\begin{tikzcd}[ampersand replacement=\&, row sep=small, column sep=small]
 \& 1 \arrow[dr] \& \\
0 \arrow[ur] \arrow[dr] \& \& 3 \\
 \& 2 \arrow[ur] \&
\end{tikzcd}
\qquad (\simeq \boldsymbol{\Delta}^1 \times \boldsymbol{\Delta}^1).
\]

\msp{page 243}
Pour toute catégorie $I$, soit
\[
\Ho_{\mathcal{M},W}(I) \text{ \ ou \ } \Ho_W(I) \overset{\text{déf}}{=} \mathcal{M}(I)W(I)^{-1},
\]
où
\[
\left\{
\begin{array}{rcl}
\mathcal{M}(I) & = & \Homu(I^\circ, \mathcal{M}) \\[1ex]
W(I) & \subseteq & \mathrm{Fl}(\mathcal{M}(I)) \text{ est formé des } u : F \longrightarrow G \text{ telles que} \\[1ex]
 & & u_i : F(i) \longrightarrow G(i) \text{ dans } W \text{ pour tout } i \text{ dans } \mathrm{Ob}\, I.
\end{array}
\right.
\]

\dgrblock{dgrb584}{One pullback with two relative meanings}
\begin{dcomment}{exposition}{p242}{One pullback with two relative meanings}
The symmetry that Grothendieck is seeking concerns the whole square. A pullback $P=A\times_S B$ has a map to $A$ and a map to $B$; a construction that remembers only its isomorphism class in a homotopy category has forgotten part of the answer. Suppose $A\to S$ is a $W$-fibration and factor $B\to S$ as $B\xrightarrow{w}\bar B\to S$, with $w\in W$ and $\bar B\to S$ a $W$-fibration. Then
\[
 A\times_S B\longrightarrow A\times_S\bar B
\]
is the pullback of $w$ along the base change of $A\to S$. It belongs to $W$. It is also a morphism over $A$, so it identifies the class of the original pullback over $A$ with the result of deriving the base-change operation on $B$. This proves the second interpretation in the text, while the first interpretation is already computed by the good representative $A\to S$. No second fibration hypothesis on $B\to S$ has been introduced.

The natural domain of this simultaneous construction is therefore the category of cospans before localization. With the manuscript's contravariant convention,
\[
 \Psi=(1\leftarrow0\to2),\qquad
 \mathcal M(\Psi)=\operatorname{Fun}(\Psi^{\mathrm{op}},\mathcal M),
\]
so an object of $\mathcal M(\Psi)$ is $A\to S\leftarrow B$. The additional vertex $3$ of $\tilde\Psi$ points, after taking the opposite, to both $1$ and $2$. A right Kan extension along $\Psi^{\mathrm{op}}\hookrightarrow\tilde\Psi^{\mathrm{op}}$ consequently puts a limit at vertex~$3$. This orientation explains the star in $i_*$: it is a right adjoint to restriction.

One should retain the distinction already emphasized in the commentary on Chapter~VI, manuscript p.~87. The category
\[
 \operatorname{Fun}(\Psi^{\mathrm{op}},\mathcal M)[W(\Psi)^{-1}]
\]
retains how the maps and their weak-equivalence comparisons occur together. It is not being replaced by $\operatorname{Fun}(\Psi^{\mathrm{op}},\mathcal M[W^{-1}])$. In the latter, a commutativity equation is only an equality of homotopy classes. The former is the diagram category on which a canonical homotopy completion can be defined. This is precisely the reason that the later derivator formalism forms diagram categories first and localizes afterwards \citep[pp.~313--315]{Groth2013}.

The coming proof accordingly has two tasks. It must show that the chosen good cospans present the localization, and then show that taking their strict pullbacks respects weak equivalences. The first task concerns all morphisms and relations in a localized category; replacing one object at a time does not settle it.
\end{dcomment}


Ainsi, le diagramme $\vcenter{\hbox{\begin{tikzcd}[ampersand replacement=\&, row sep=small, column sep=small]
 \& X_1 \arrow[dl] \\
S \& \\
 \& X_2 \arrow[ul]
\end{tikzcd}}}$ définit un objet $\xi$ de $\mathcal{M}(\Psi)$, donc de $\Ho_W(\Psi)$, et le diagramme $\vcenter{\hbox{\begin{tikzcd}[ampersand replacement=\&, row sep=small, column sep=small]
 \& X_1 \arrow[dl] \& \\
S \& \& Y \arrow[ul] \arrow[dl] \\
 \& X_2 \arrow[ul] \&
\end{tikzcd}}}$ peut de même s'interpréter comme un objet $\tilde{\xi}$ de $\Ho_W(\tilde{\Psi})$, dont la restriction à $\Psi$ est $\xi$.

Ceci posé, j'aimerais définir un foncteur
\[
i_* : \Ho_W(\Psi) \longrightarrow \Ho_W(\tilde{\Psi})
\]
(où
\[
i : \Psi \longrightarrow \tilde{\Psi}
\]
est l'inclusion, qui est une immersion ouverte), lequel doit correspondre, dans le cas d'un objet $\vcenter{\hbox{\begin{tikzcd}[ampersand replacement=\&, row sep=small, column sep=small]
 \& X_1 \arrow[dl] \\
S \& \\
 \& X_2 \arrow[ul]
\end{tikzcd}}}$ de $\Ho_W(\Psi)$ avec $p_1 \in \mathrm{Fib}_W$ \texttt{[}$:$ $X_1 \longrightarrow S$\texttt{]}, à l'opération de produit fibré (\footnote[131]{\textbf{NB} Cette construction, qui se fait ici assez lourdement, devient immédiate, et aussi ses principales propriétés, si on admet que la catégorie à fibrations $(\mathcal{M}, W, \mathrm{Fib}_W)$ définit un \textit{dérivateur} sur le domaine des ensembles ordonnés finis, ce qui se prouve en effet assez simplement (avec pas plus de mal que je ne me donne ici).}).

\printedpage{158}
% head: A. Grothendieck : Dérivateurs, Ch. XII, Édition des Universités de Montpellier II et Paris VII
% foot: Version 30 mars 2003

\msp{page 244}
La construction de ce foncteur $i_*$ devrait être essentiellement sorital, avec ce dont nous disposons. Nous allons définir d'abord, pour $S$ \textit{fixé,} un foncteur canonique
\[
\Ho^{\lc}_W(S) \times \Ho^{\lc}_W(S) \longrightarrow \Ho_{W,S}(\Psi), \tag*{$(*)$}
\]
où le deuxième membre est défini, par analogie avec les $\Ho_W(I)$, par
\[
\Ho_{W,S}(I) \overset{\text{déf}}{=} ((\mathcal{M}/S)(I))(W_S(I))^{-1}
\]
- i.e. c'est la même définition, sauf que $(\mathcal{M}, W)$ est remplacé par $(\mathcal{M}/S, W_S)$. Quant à $\Phi = \vcenter{\hbox{\begin{tikzcd}[ampersand replacement=\&, row sep=small, column sep=small]
1 \arrow[dr] \& \\
 \& 3 \\
2 \arrow[ur] \&
\end{tikzcd}}}$ , on le considère comme une sous-catégorie fermée de $\tilde{\Psi}$, qui s'en déduit en lui rajoutant un objet initial 0 (tout comme $\tilde{\Psi}$ se déduit de $\Psi$ en lui rajoutant un objet final 3). D'autre part, on a un foncteur tautologique
\[
\Ho_{W,S}(\Phi) \longrightarrow \Ho_W(\tilde{\Psi}), \tag*{$(*')$}
\]
qui sur les objet consiste à compléter un diagramme $\vcenter{\hbox{\begin{tikzcd}[ampersand replacement=\&, row sep=small, column sep=small] % [sic] les objet -> les objets, plural agreement
X_1 \& \\
 \& Y \arrow[ul] \arrow[dl] \\
X_2 \&
\end{tikzcd}}}$ de $\mathcal{M}/S$ en le diagramme carré $\vcenter{\hbox{\begin{tikzcd}[ampersand replacement=\&, row sep=small, column sep=small]
 \& X_1 \arrow[dl] \& \\
S \& \& Y \arrow[ul] \arrow[dl] \\
 \& X_2 \arrow[ul] \&
\end{tikzcd}}}$ de $\mathcal{M}$.

\msp{page 245}
Le composé de $(*)$ et de $(*')$ donnera donc un foncteur
\[
\Ho^{\lc}_W(S) \times \Ho^{\lc}_W(S) \longrightarrow \Ho_W(\tilde{\Psi}), \tag*{$(**)$}
\]
très proche déjà du foncteur $i_*$ à construire. Pour construire $(*)$, on interprète (à équivalence de catégories près) le premier membre comme
\[
(\Fibu_W S \times \Fibu_W S)(W_S \times W_S)^{-1},
\]
donc il faut construire un foncteur
\[
\Fibu_W S \times \Fibu_W S \longrightarrow \Ho_{W,S}(\Phi),
\]
qui transforme flèches dans $W_S \times W_S$ en isomorphismes. On prendra le foncteur
\[
(X_1, X_2) \ \longmapsto \
\left(
\begin{tikzcd}[ampersand replacement=\&]
X_1 \& X_1 \times_S X_2 \arrow[l] \arrow[d] \\
 \& X_2
\end{tikzcd}
\right),
\]
\printedpage{159}
% head: A. Grothendieck : Dérivateurs, Ch. XII, Édition des Universités de Montpellier II et Paris VII
% foot: Version 30 mars 2003
et il faut seulement vérifier que si
\[
X_1 \xrightarrow{\alpha_1} X_1', \qquad X_2 \xrightarrow{\alpha_2} X_2'
\]
sont deux flèches dans $\Fibu_W S$, d'où
\[
X_1 \times_S X_2 \xrightarrow{\alpha = \alpha_1 \times_S \alpha_2} X_1' \times_S X_2',
\]
alors
\[
\alpha_1 \in W,\ \alpha_2 \in W \qquad \Longrightarrow \qquad \alpha \in W,
\]
ce qui est immédiat sur la factorisation
\[
X_1 \times_S X_2 \quad \underbrace{\longrightarrow}_{\substack{\in W \\ \text{car } X_2 \in \Fibu_W S}} \quad X_1' \times_S X_2 \quad \underbrace{\longrightarrow}_{\substack{\in W \\ \text{car } X_1' \in \Fibu_W S}} \quad X_1' \times_S X_2'.
\]

\msp{page 246}
On a prouvé la

\textbf{Proposition 18.} \textit{Soit $S$ dans $\mathcal{M}$. Il existe, à isomorphisme unique près, un foncteur et un seul}
\[
\Ho^{\lc}_W(S) \times \Ho^{\lc}_W(S) \xrightarrow{\pi} \Ho_{W,S}(\Phi)
\]
\textit{donnant lieu à un carré commutatif à isomorphisme près}
\[
\resizebox{\linewidth}{!}{$\begin{tikzcd}[ampersand replacement=\&, column sep=large]
 \& {\scriptstyle (X_1,X_2)\ \longmapsto\ \begin{array}{@{}c@{}c@{}c@{}} \scriptstyle X_1 & & \\ & \nwarrow & \\ & & \scriptstyle X_1 \times_S X_2 \\ & \swarrow & \\ \scriptstyle X_2 & & \end{array}} \& \\[-3ex]
\Fibu_W S \times \Fibu_W S \arrow[rr] \arrow[d] \& \& (\mathcal{M}/S)(\Phi) \arrow[d] \\
\Ho^{\lc}_W(S) \times \Ho^{\lc}_W(S) \arrow[rr] \& \& \Ho_{W,S}(\Phi) \\[-4ex]
(= (\mathcal{M}/S)(W_S)^{-1} \times (\mathcal{M}/S)(W_S)^{-1}) \& \& (= \mathcal{M}/S(\Phi)(W_S(\Phi))^{-1}).
\end{tikzcd}$}
\]
\textit{Pour $X_1$, $X_2$ dans $\mathcal{M}/S$, le diagramme}
\[
\pi(X_1, X_2) \overset{\text{déf}}{=} \pi(\ho^{\lc}_{W,S}(X_1), \ho^{\lc}_{W,S}(X_2))
\]
\textit{se `calcule' ainsi, à isomorphisme canonique près : on choisit}
\[
X_1 \xrightarrow{\alpha_1} \bar{X}_1, \qquad X_2 \xrightarrow{\alpha_2} \bar{X}_2 \qquad \textit{flèches dans } \mathcal{M}/S
\]
\textit{avec}
\[
\alpha_1, \alpha_2 \in W_S, \qquad \bar{X}_1, \bar{X}_2 \in \mathrm{Ob}\, \Fibu_W S,
\]
\textit{ce qui définit}
\[
\pi(\alpha_1, \alpha_2) : \pi(X_1, X_2) \xrightarrow{\sim} \pi(\bar{X}_1, \bar{X}_2) = \ho_{S,W,\Phi}
\left(
\begin{tikzcd}[ampersand replacement=\&]
\bar{X}_1 \& \bar{X}_1 \times_S \bar{X}_2 \arrow[l] \arrow[d] \\
 \& \bar{X}_2
\end{tikzcd}
\right).
\]

% ---- from Der-ChXIIc-30-03-03_p046-050.tex
\setcounter{footnote}{131}

\printedpage{160}
% head: A. Grothendieck : Dérivateurs, Ch. XII, Édition des Universités de Montpellier II et Paris VII
% foot: Version 30 mars 2003

\msp{page 247}

\dgrblock{dgrb585}{A product together with its projections}
\begin{dcomment}{exposition}{p246}{A product together with its projections}
The intermediate construction over a fixed base is more than a calculation of the object $A\times_S B$. Grothendieck keeps both projections in the value of $\pi$. Thus $\pi$ lands in the localized category of $\Phi^{\mathrm{op}}$-diagrams in $\mathcal M/S$, where $\Phi=(1\to3\leftarrow2)$. The occurrence of $\Psi$ in the first display of ms.~244 is an edition slip: the next functor, and Proposition~18 itself, require $\Phi$.

Write $\mathcal C_S$ for the full subcategory of $\mathcal M/S$ on the $W$-fibrations. If $a:A\to A'$ and $b:B\to B'$ are weak equivalences in this category, the comparison of products has the factorization
\[
 A\times_S B\xrightarrow{a\times_S1}
 A'\times_S B\xrightarrow{1\times_S b}
 A'\times_S B'.
\]
The first arrow is a pullback of $a$ along a base change of $B\to S$; the second is a pullback of $b$ along a base change of $A'\to S$. Both belong to $W$. The same morphism of spans has $a$ and $b$ at its other vertices, so it is an objectwise weak equivalence. This proves exactly the compatibility needed to descend the entire span functor.

There is also a small categorical point behind the source's product of localizations. For relative categories $(\mathcal C,U)$ and $(\mathcal D,V)$, with identities in $U,V$, there is a canonical equivalence
\[
 (\mathcal C\times\mathcal D)[(U\times V)^{-1}]
 \simeq\mathcal C[U^{-1}]\times\mathcal D[V^{-1}].
\]
Indeed a functor out of $\mathcal C\times\mathcal D$ inverts $U\times V$ precisely when it inverts $(u,1)$ and $(1,v)$ separately. Localizing successively in the two variables gives the displayed universal property, including natural transformations. Applied to $\mathcal C_S$ twice, this justifies the domain used in Proposition~18. The localization equivalence between $\mathcal C_S$ and the relative homotopy category is an earlier input to this paragraph, not a consequence of the displayed product calculation.

The phrase ``unique up to unique isomorphism'' is understood relative to the specified comparison with the functor on good representatives. Without that comparison, a functor can have automorphisms and the uniqueness assertion would be too strong. With it, full faithfulness of restriction along a localization supplies the required uniqueness. A pair of replacement arrows $A\to\bar A$, $B\to\bar B$ therefore gives the canonical comparison written by Grothendieck; it does not require a preferred strict map between every two independently chosen replacements.

This proposition is numbered~18 a second time in the edition. We distinguish it from the pasting proposition on ms.~236 by its page and by the name ``fixed-base product proposition''. Its content will also be used when the base itself varies, where a further argument is needed.
\end{dcomment}


\textbf{Corollaire.} {\itshape Description du foncteur
\[
(\xi_1,\xi_2)\longmapsto\xi_1\times_S\xi_2 : \Ho^{\lc}_W(S)\times\Ho^{\lc}_W(S)\longrightarrow\Ho^{\lc}_W(S).
\]
Ce foncteur est bel et bien un produit cartésien dans} $\Ho^{\lc}_W(S)$ (\footnote{mais c'est un cas particulier des cas de $\Ho_W$ (si $\mathcal{M}$ a des produits binaires), en remplaçant $\mathcal{M}$ par $\mathcal{M}/S$.}), {\itshape où les produits finis existent} (\footnote{Dire aussi que les foncteurs $f^* : \Ho^{\lc}_W(S)\longrightarrow\Ho^{\lc}_W(S')$ (pour $f : S'\longrightarrow S$) commutent aux produits finis.}).

Pour prouver ce dernier point, posant
\[
\mathcal{H}=\Ho^{\lc}_W(S),
\]
il faut montrer que les foncteurs
\[
\begin{tikzcd}[ampersand replacement=\&]
\mathcal{H} \arrow[r, shift left, "\delta"] \& \mathcal{H}\times\mathcal{H} \arrow[l, shift left, "\pi"]
\end{tikzcd}
\]
sont adjoints, avec
\[
\begin{array}{rclcrcl}
\delta & = & \mathrm{diag}_{\mathcal{H}} & : & \xi & \longmapsto & (\xi,\xi)\\
\pi & = & \text{`produit'} & : & (\xi_1,\xi_2) & \longmapsto & \xi_1\times_S\xi_2.
\end{array}
\]
On va définir les homomorphismes d'adjonction
\[
\begin{array}{rcllrcl}
\mathrm{id}_{\mathcal{H}} & \xrightarrow{\alpha} & \pi\delta & \text{i.e.} & \xi & \xrightarrow{\alpha_\xi} & \xi\times_S\xi\\
\mathrm{id}_{\mathcal{H}\times\mathcal{H}} & \xleftarrow{\beta} & \delta\pi & \text{i.e.} & (\xi_1,\xi_2) & \xleftarrow{\beta_{\xi_1,\xi_2}} & (\xi_1\times_S\xi_2,\xi_1\times_S\xi_2).
\end{array}
\]
Pour ceci, on décrit $\mathcal{H}$ comme
\[
\mathcal{H}=(\uFib_W S)(W_S)^{-1}
\]
et on note que $\delta$, $\pi$ proviennent de foncteurs adjoints au niveau de $\mathcal{H}_0=\uFib_W S$
\[
\begin{tikzcd}[ampersand replacement=\&]
\mathcal{H}_0 \arrow[r, shift left, "\delta_0"] \& \mathcal{H}_0\times\mathcal{H}_0 \arrow[l, shift left, "\pi_0"]
\end{tikzcd},
\]
lesquels donnent lieu à $\alpha_0$, $\beta_0$
\[
\begin{array}{rcl}
\mathrm{id}_{\mathcal{H}_0} & \xrightarrow{\alpha_0} & \pi_0\delta_0\\
\mathrm{id}_{\mathcal{H}_0\times\mathcal{H}_0} & \xleftarrow{\beta_0} & \delta_0\pi_0.
\end{array}
\]

\printedpage{161}
% head: A. Grothendieck : Dérivateurs, Ch. XII, Édition des Universités de Montpellier II et Paris VII
% foot: Version 30 mars 2003

\msp{page 248}

On définit $\alpha$, $\beta$ par passage aux catégories de fractions pour $\alpha_0$, $\beta_0$. Les 2 compatibilités, étant vérifiées pour $\alpha_0$, $\beta_0$, le sont aussi pour $\alpha$, $\beta$. Q.e.d.

Je dégage ici le

\textbf{Lemme.} {\itshape Soit \begin{tikzcd}[ampersand replacement=\&]
\mathcal{X} \arrow[r, shift left, "f"] \& \mathcal{Y} \arrow[l, shift left, "g"]
\end{tikzcd} un couple de foncteurs adjoint, et soient $\Sigma$, $\Sigma'$ des localiseurs dans $\mathcal{X}$, $\mathcal{Y}$ tels que $f(\Sigma)\subseteq\Sigma'$, $g(\Sigma')\subseteq\Sigma$, de sorte qu'on a un couple de foncteurs
\[
\begin{tikzcd}[ampersand replacement=\&]
\mathcal{X}\Sigma^{-1} \arrow[r, shift left, "\bar{f}"] \& \mathcal{Y}\Sigma'^{-1} \arrow[l, shift left, "\bar{g}"]
\end{tikzcd}.
\]
Alors $\bar{f}$, $\bar{g}$ sont adjoints, les flèches d'adjonction étant déduites de celles pour $f$, $g$.}

\textbf{Corollaire.} {\itshape Soit $(\mathcal{M},W)$ une catégorie de modèles, supposons que dans $\mathcal{M}$ les produits binaires (resp. les produits finis) existent, et que pour $u_1 : X_1\longrightarrow X'_1$ et $u_2 : X_2\longrightarrow X'_2$, si $u_1,u_2\in W$, alors $u_1\times u_2 : X_1\times X_2\longrightarrow X'_1\times X'_2$ est dans $W$. Alors dans $\mathcal{M}W^{-1}$ les produits binaires (resp. produits finis) existent, et le foncteur localisation $\mathcal{M}\longrightarrow\mathcal{M}W^{-1}$ y commute.}

On applique le lemme à
\[
\begin{tikzcd}[ampersand replacement=\&]
\mathcal{M} \arrow[r, shift left, "\delta"] \& \mathcal{M}\times\mathcal{M} \arrow[l, shift left, "\pi"]
\end{tikzcd}
\qquad \text{et } W,\ W\times W,
\]
et pour l'objet final,
\[
\begin{tikzcd}[ampersand replacement=\&]
\mathcal{M} \arrow[r, shift left, "e"] \& e \arrow[l, shift left, "\bullet"]
\end{tikzcd}
\qquad\qquad \text{et } W,\ W_0=\{\mathrm{id}_e\}.
\]

\msp{page 249}

\dgrblock{dgrb586}{Why the product remains a categorical product}
\begin{dcomment}{exposition}{p248}{Why the product remains a categorical product}
The assertion about products in $\Ho_W^{\mathrm{lc}}(S)$ is stronger than the existence of a homotopically well-defined binary operation. Grothendieck proves the universal property by descending an adjunction. This is an efficient place to make the small localization argument completely explicit.

Suppose $F:\mathcal C\rightleftarrows\mathcal D:G$ is an adjunction, with $F$ left adjoint to $G$, and both functors preserve the chosen weak equivalences. Its unit $\eta:1\to GF$ and counit $\epsilon:FG\to1$ have images in the localizations. Their naturality for an inverted weak equivalence follows from their naturality for that weak equivalence: multiply the naturality square by the inverse image of the equivalence. Naturality then holds for any composite of ordinary arrows and formal inverses, which are the morphisms of the localization. The identities
\[
 \epsilon_{F X}F(\eta_X)=1_{F X},\qquad
 G(\epsilon_Y)\eta_{G Y}=1_{G Y}
\]
are equations already true before localization, hence remain true afterwards. The descended unit and counit establish $\bar F\dashv\bar G$. No lifting property and no fullness assertion about the localization functor are needed.

Apply this argument to the diagonal functor on $\mathcal C_S$ and its right adjoint, the ordinary product over $S$. The preceding block proves that this right adjoint preserves weak equivalences. The unit is the diagonal $A\to A\times_S A$; the counit consists of the two projections. Consequently the operation constructed from good representatives is the categorical product in the localized relative category. Applying the same argument to the unique functor $\mathcal C_S\to e$ and the functor selecting $S\to S$ proves that this object remains terminal. This supplies all finite products.

The statement in footnote~133 has an equally concrete proof. Pullback along $f:S'\to S$ sends good objects to good objects and, by the preceding relative invariance theorem, sends weak equivalences between them to weak equivalences. The canonical strict isomorphism
\[
 f^*(A\times_S B)\cong f^*A\times_{S'}f^*B
\]
is compatible with both projections; also $f^*(S\to S)=(S'\to S')$. Their images after localization are therefore the comparison isomorphisms for a functor preserving finite products. This proof concerns the product in the relative homotopy category. It does not assert that an arbitrary strict pullback in $\mathcal M$ is already a homotopy pullback.
\end{dcomment}


Je vais maintenant construire le foncteur
\[
i_* : \Ho_W(\Psi)\longrightarrow\Ho_W(\tilde{\Psi}).
\]
Il me faut d'abord prouver le

\textbf{Lemme 1.} {\itshape Soit $\mathcal{F}$ la sous-catégorie pleine de $\mathcal{M}(\Psi)$, catégorie des diagrammes \begin{tikzcd}[ampersand replacement=\&, row sep=small, column sep=small]
 \& X_1 \arrow[dl, "p_1"'] \\
S \& \\
 \& X_2 \arrow[ul, "p_2"]
\end{tikzcd} dans $S$ \texttt{[plutôt dans }$\mathcal{M}$\texttt{, sur }$S$\texttt{]}, formée de ceux pour lesquels $p_1$, $p_2$ sont dans Fib$_W$. Soit $W_{\mathcal{F}}$ le localiseur induit sur $\mathcal{F}$ par $W(\Psi)$ ($W$-équivalences terme à terme). Alors
\[
\mathcal{F}W_{\mathcal{F}}^{-1}\xrightarrow{\sim}\mathcal{M}(\Psi)W(\Psi)^{-1}\ (\overset{\text{déf}}{=}\Ho_W(\Psi)).
\]}

\printedpage{162}
% head: A. Grothendieck : Dérivateurs, Ch. XII, Édition des Universités de Montpellier II et Paris VII
% foot: Version 30 mars 2003

Nous allons prouver ce lemme plus bas. En l'admettant, pour construire $i_*$, il suffit de construire
\[
\mathcal{F}\longrightarrow\Ho_W(\tilde{\Psi})
\]
qui applique $W_{\mathcal{F}}$ en des isomorphismes. Je prendrai, bien sûr,
\[
\begin{gathered}
\begin{tikzcd}[ampersand replacement=\&, row sep=small, column sep=small]
 \& X_1 \arrow[dl] \& \& \& \& X_1 \arrow[dl] \& \\
S \& \& \longmapsto \& \& S \& \& X_1\times_S X_2 \arrow[ul] \arrow[dl] \\
 \& X_2 \arrow[ul] \& \& \& \& X_2 \arrow[ul] \&
\end{tikzcd}\\
\mathcal{F}\quad\longrightarrow\quad\Ho_W(\Psi).
\end{gathered}
\] % [sic] \Ho_W(\Psi) -> \Ho_W(\tilde{\Psi}), the image is a diagram of type \tilde{\Psi} and the functor to be constructed is \mathcal{F}\longrightarrow\Ho_W(\tilde{\Psi})
Vérifions qu'il transforme flèches de $W_{\mathcal{F}}$ en isomorphismes.

\msp{page 250}

\dgrblock{dgrb587}{The deferred lemmas with functorial factorization}
\begin{dcomment}{exposition}{p249}{The deferred lemmas with functorial factorization}
Grothendieck has deliberately separated the existence of useful representatives from the assertion that they compute the whole localization. Lemmas~1 and~2 assert the latter. Their proofs are deferred, and the promise on ms.~257 does not complete them in the pages supplied here. There is a particularly short completion when the factorization axiom is available functorially. We state that extra hypothesis explicitly, since the axiom on ms.~240 does not include it.

Let $D=\mathcal M(\Psi)$, and suppose that the assignment
\[
 (A\xrightarrow p S)\longmapsto
 (A\xrightarrow{j_p}R(p)\xrightarrow{r_p}S)
\]
is a functor on the arrow category, with $j_p\in W$, $r_p\in\mathrm{Fib}_W$, and $r_pj_p=p$. For $d=(A\xrightarrow p S\xleftarrow q B)$, set
\[
 R_Dd=(R(p)\xrightarrow{r_p}S\xleftarrow{r_q}R(q)).
\]
This gives a functor $R_D:D\to\mathcal F$ and a natural transformation $j:1_D\to J R_D$, where $J:\mathcal F\hookrightarrow D$ is inclusion. Every component of $j$ is an objectwise weak equivalence. If $a:d\to d'$ is in $W(\Psi)$, naturality gives
\[
 (J R_Da)j_d=j_{d'}a.
\]
At each vertex the two maps $j_d,j_{d'}$ and the map $a$ belong to $W$; two-out-of-three therefore implies that $R_Da$ is in $W_{\mathcal F}$. The functor $R_D$ consequently descends to localization. Restricted to $\mathcal F$, the same natural transformation has both source and target in $\mathcal F$. After localization it supplies an isomorphism $1\cong \bar R_D\bar J$, while its unrestricted version supplies $1\cong\bar J\bar R_D$. These are precisely the two isomorphisms that show $\bar J$ is an equivalence. This proves Lemma~1 with the additional hypothesis. Over a fixed base $S$, the same argument applies to every full intermediate category between $\mathrm{Fib}_W(S)$ and $\mathcal M/S$: the replacement always lands in $\mathrm{Fib}_W(S)$ and the natural comparison stays in the intermediate category. It therefore proves the strengthened replacement assertion~(c) recalled on ms.~240 in this functorial setting.

For Lemma~2 use the identical construction on the three vertices of a square belonging to $\Psi$. Keep its fourth vertex $T$ unchanged, and replace its maps to $A$ and $B$ by their composites with $j_p$ and $j_q$. Their composites to $S$ still agree, so this is a strictly commutative square. It defines a functor $R_E:\mathcal M(\tilde\Psi)\to\mathcal G$ with a natural objectwise weak equivalence $1\to K R_E$, whose component at $T$ is the identity. The same two-out-of-three argument makes $R_E$ homotopical, and its restriction to $\mathcal G$ gives the other inverse comparison. Hence $\mathcal G[W_{\mathcal G}^{-1}]\to\Ho_W(\tilde\Psi)$ is an equivalence.

This argument also specifies what the arbitrary choice of a factorization would fail to provide: not the value on an object, but the functor and the natural transformation used in the proof. It uses neither an assumption that every object is fibrant nor a saturation theorem identifying $W$ with all maps inverted by localization. The required input is a functorial factorization, which is a separate sufficient hypothesis.
\end{dcomment}


\dgrblock{dgrb588}{A nonfunctorial repair by derivability}
\begin{dcomment}{later}{p249}{A nonfunctorial repair by derivability}
The passage to diagrams can be made without functorial choices, and the later theory also removes the need to add global fibrancy to the written hypotheses. The appropriate structure is a left derivable category in the sense of \citet[Definition~1.1]{Cisinski2010}. We check its one potentially hidden axiom. Let $p:X\to Y$ belong to $W\cap\mathrm{Fib}_W$, and let $u:Y'\to Y$ be arbitrary. Factor $u$ as $Y'\xrightarrow{j}Z\xrightarrow{r}Y$, where $j\in W$ and $r\in\mathrm{Fib}_W$. Since $r^*$ preserves $W$, the map $X\times_Y Z\to Z$ is in $W$. The map $X\times_Y Y'\to X\times_Y Z$ is the pullback of $j$ along a base change of $p$, hence is also in $W$. The commutative square with bottom map $j$ and two-out-of-three now give $X\times_Y Y'\to Y'\in W$. Thus trivial fibrations remain trivial under arbitrary pullback.

The other derivability axioms are immediate from the standing assumptions: a terminal object exists, weak equivalences contain isomorphisms and satisfy two-out-of-three, fibrations compose and are stable under the available pullbacks, and the required factorization is the axiom of ms.~240. These properties are stronger than the restricted pullback and factorization requirements of Definition~1.1. The precise later results are the finite-diagram structures and localization comparison of Theorem~1.30 and Corollaries~1.31--1.32, followed by the right-Kan-extension theorem, Proposition~2.6; Theorem~2.21 packages the construction as a one-sided derivator on finite direct categories.

This also settles the strengthened fixed-base condition~(c) recalled from ms.~218. The slice $\mathcal M/S$, with the inherited fibrations and weak equivalences, is left derivable and has $\mathrm{Fib}_W(S)$ as its full subcategory of fibrant objects. Every full intermediate category $\mathcal E$ containing these objects is left derivable with the same fibrant objects: the pullbacks demanded by the axioms take place among fibrant objects, and the middle object in a factorization with fibrant target is again fibrant. Proposition~1.8 identifies the localization of each such category with that of its fibrant objects. It follows that both inclusions in condition~(c) are localization equivalences. Condition~(b), the equality of ordinary and universal weak equivalences between good objects, is the relative invariance argument dual to Proposition~16; together they give~(a).

For clarity, the particular inclusions $\mathcal F\subset D$ and $\mathcal G\subset E$ can be checked by that paper's approximation theorem, Theorem~3.12. We give the verification, since $\mathcal G$ imposes no fibration condition on its fourth vertex. On $\mathcal F$ take a morphism $(a_0,a_1,a_2)$ to be a fibration when the three maps
\[
 S\to S',\qquad
 A\to S\times_{S'}A',\qquad
 B\to S\times_{S'}B'
\]
are fibrations. On $\mathcal G$ use these same conditions and additionally require the map $T\to T'$ between fourth vertices to be a fibration. Take weak equivalences objectwise in both cases. An object of $\mathcal F$ is fibrant precisely when its base $S$ is fibrant in $\mathcal M$. An object of $\mathcal G$ is fibrant precisely when both $S$ and its fourth vertex $T$ are fibrant. Its remaining vertices are then fibrant because their maps to $S$ are fibrations.

Here is the factorization needed both for these structures and for approximation. Given a map from an arbitrary cospan to a fibrant good cospan $(A'\to S'\leftarrow B')$, first factor its base map as
\[
 S\xrightarrow{\sim}S_1\twoheadrightarrow S'.
\]
Factor the resulting maps $A\to S_1\times_{S'}A'$ and $B\to S_1\times_{S'}B'$ as weak equivalences followed by fibrations. Call the middle objects $A_1,B_1$. Their maps to $S_1$ are fibrations, and the resulting cospan maps to the target by a fibration of the displayed kind. This constructs an objectwise weak equivalence from the original cospan to an object of $\mathcal F$, compatibly with the given map to the target.

For a map of squares with fibrant target in $\mathcal G$, make those three factorizations and put
\[
 P_1=A_1\times_{S_1}B_1,\qquad P'=A'\times_{S'}B'.
\]
The map $P_1\to P'$ is a fibration. Indeed it factors through $S_1\times_{S'}P'$; the first map is a composite of pullbacks of the two relative leg fibrations, and the second is a pullback of $S_1\to S'$. The source corner $T$ maps canonically to $P_1\times_{P'}T'$. Factor this map into $T\xrightarrow{\sim}T_1$ followed by a fibration. Since $P_1\times_{P'}T'\to T'$ is itself a fibration, so is $T_1\to T'$. The middle square lies in $\mathcal G$ and is fibrant: its base maps fibrantly to the fibrant target base, and $T_1$ maps fibrantly to the fibrant object $T'$. Thus the factorization has all the required properties.

The remaining derivability axioms can be checked on these same relative maps. Composition is a composite of a relative map and a pullback of the next one. A pullback of a fibration is computed vertex by vertex; at either leg its relative map is the corresponding pullback of the original relative map, so the resulting object remains good. For a trivial fibration, the relative leg maps are also weak equivalences, by two-out-of-three applied to the weak base map and the good target leg. Their pullbacks remain weak equivalences, as do the fourth-vertex maps. This verifies the trivial-fibration axiom. Factoring a diagonal of a fibrant object by the construction above supplies its path object. The ambient categories of cospans and squares carry the objectwise fibration structures of Cisinski's Corollary~1.31; the nontrivial factorization axiom there is supplied by the finite inverse diagram construction. The inclusions preserve fibrations, trivial fibrations, terminal objects and their pullbacks; they are exact and reflect weak equivalences.

Finally, Cisinski's approximation condition AP2 asks, for a map from an ambient fibrant object to an included fibrant object, for a weak-equivalence comparison with an object of the smaller category over the same target. The explicit factorizations just constructed supply it, with the preliminary weak equivalence of the ambient source taken to be its identity. AP1 is reflection of weak equivalences. Theorem~3.12 now proves both localization equivalences. This proves the two deferred lemmas under the written factorization and pullback hypotheses. The derivable-category argument controls the morphisms without either global fibrancy or functorial choices of replacements.
\end{dcomment}


\[
\begin{tikzcd}[ampersand replacement=\&, row sep=small, column sep=small]
 \& X_1\times_S X_2 \arrow[dl] \arrow[dd] \& \& X'_1\times_{S'}X'_2 \arrow[ll, "f_3"'] \arrow[dl] \arrow[dd] \\
X_1 \arrow[dd] \& \& X'_1 \arrow[ll, crossing over, "f_1"'] \& \\
 \& X_2 \arrow[dl] \& \& X'_2 \arrow[ll, "f_2"'] \arrow[dl] \\
S \& \& S' \arrow[ll, "f_0"'] \arrow[from=uu, crossing over] \&
\end{tikzcd}
\]

Il faut prouver que si $f_0,f_1,f_2\in W$, alors $f_3\in W$. Mais considérons l'image inverse du carré cartésien $Q$ par $f_0 : S'\longrightarrow S$, d'où un carré cartésien $Q_{S'}$ et une factorisation
\[
\begin{tikzcd}[ampersand replacement=\&]
Q \& Q_{S'} \arrow[l, "\tilde{f}"'] \& Q' \arrow[l, "g"'] \arrow[ll, bend left, "f"]
\end{tikzcd}
\]
de $Q'\longrightarrow Q$. Comme $f_0\in W$, et que $X_1$, $X_2$, $X_1\times_S X_2$ sont dans $\uFib_W S$, les flèches horizontales réalisant $Q_{S'}\longrightarrow Q$ sont toutes dans $W$. Ceci, plus le fait que $f_1,f_2\in W$, impliquent que dans le morphisme $g$, les flèches $g_1$, $g_2$ sont dans $W$. Il en résulte que la flèche
\[
g_3 : X'_1\times_{S'}X'_2\longrightarrow(X_1\times_S X_2)\times_S S'
\]
l'est aussi (théorie du produit dans $\uFib_W S'$ et $\Ho^{\lc}_W(S)$), donc $f_3=\tilde{f}_3g_3\in W$, q.e.d.

\msp{page 251}

\dgrblock{dgrb589}{The base is allowed to move}
\begin{dcomment}{exposition}{p250}{The base is allowed to move}
Having obtained good representatives, Grothendieck must show that their pullbacks vary homotopically in the base as well as in the two inputs. This is the extra content of the argument on ms.~250. The fixed-base product proposition cannot by itself be applied to a morphism of cospans whose middle objects differ.

Suppose a morphism of good cospans has components
\[
 a:A'\to A,\qquad b:B'\to B,\qquad s:S'\to S
\]
in $W$. Pull the target cospan back along $s$, obtaining
\[
 \tilde A=S'\times_S A,\quad
 \tilde B=S'\times_S B,\quad
 \tilde P=S'\times_S(A\times_S B).
\]
The projections $\tilde A\to A$ and $\tilde B\to B$ are weak equivalences because $A\to S$ and $B\to S$ are $W$-fibrations. The map $A\times_S B\to S$ is a composite of $W$-fibrations, so $\tilde P\to A\times_S B$ is also in $W$. The given maps factor through $\tilde A$ and $\tilde B$; two-out-of-three shows that $A'\to\tilde A$ and $B'\to\tilde B$ lie in $W$. These are now maps over the same base $S'$. The product calculation gives
\[
 A'\times_{S'}B'\longrightarrow
 \tilde A\times_{S'}\tilde B\cong\tilde P
 \quad\text{in }W.
\]
Composition proves that the original map of pullbacks is a weak equivalence. All four components of the completed square have therefore been checked. The final parenthesis of the source should refer to $\Ho_W^{\mathrm{lc}}(S')$, since this intermediate product is taken over $S'$; the printed $S$ is retained as an edition slip.

The calculation explains why one can descend the strict completion functor on $\mathcal F$. It also shows that its homotopicality is a theorem about a whole diagram, not an assertion that the pullback functor preserves arbitrary weak equivalences in each uncontrolled variable. The good-map assumptions are used at the three projections above and at the two fixed-base product maps. In Brown's setting the corresponding pullback invariance is already the content of \citet[Part~I, \S4, Lemmas~1--2, pp.~428--429]{Brown1973}.
\end{dcomment}


Notons qu'on a un isomorphisme canonique
\[
\alpha : i^*i_*\xrightarrow{\sim}\mathrm{id}_{\Ho_W(\Phi)}. \tag*{$(*)$}
\]
J'ai envie de prouver que les foncteurs
\[
\begin{tikzcd}[ampersand replacement=\&]
\Ho_W(\Phi) \arrow[r, shift right, "i_*"'] \& \Ho_W(\tilde{\Phi}) \arrow[l, shift right, "i^*"']
\end{tikzcd}
\]
\printedpage{163}
% head: A. Grothendieck : Dérivateurs, Ch. XII, Édition des Universités de Montpellier II et Paris VII
% foot: Version 30 mars 2003
sont adjoints, et que $(*)$ est une des flèches d'adjonction. Il faut définir l'autre flèche d'adjonction
\[
\beta : i_*i^*\longleftarrow\mathrm{id}_{\Ho_W(\tilde{\Phi})}. \tag*{$(**)$}
\]
Mais pour ceci, je dois utiliser le

\textbf{Lemme 2.} Soit $\mathcal{G}$ la sous-catégorie pleine de $\mathcal{M}(\tilde{\Phi})$, formée des \begin{tikzcd}[ampersand replacement=\&]
X_1 \arrow[d, "p_1"'] \& X_3 \arrow[l, "q_2"'] \arrow[d, "q_1"] \\
S \& X_2 \arrow[l, "p_2"]
\end{tikzcd} tels que $p_1,p_2\in\mathrm{Fib}_W$ et soit $W_{\mathcal{G}}$ le localiseur induit par $W(\tilde{\Phi})$ ($W$-équivalences terme à terme). Alors l'inclusion $\mathcal{G}\longrightarrow\mathcal{M}(\tilde{\Phi})$ induit une équivalence
\[
\mathcal{G}W_{\mathcal{G}}^{-1}\xrightarrow{\sim}\mathcal{M}(\tilde{\Phi})W(\tilde{\Phi})^{-1}\overset{\text{déf}}{=}\Ho_W(\tilde{\Phi}).
\]

Je vais pour le moment admettre ce

\msp{page 252}

lemme, tout comme le lemme 1 (p. 249). Je vais considérer $i^*$, $i_*$ comme des flèches entre les categories $\mathcal{F}W_{\mathcal{F}}^{-1}$, $\mathcal{G}W_{\mathcal{G}}^{-1}$. Elles sont induites par les flèches % [sic] categories -> catégories, missing accent
\[
\begin{tikzcd}[ampersand replacement=\&]
\mathcal{F} \arrow[r, shift right, "i_{0\,*}"'] \& \mathcal{G} \arrow[l, shift right, "i_0^*"']
\end{tikzcd}
\]
entre $\mathcal{F}$ et $\mathcal{G}$ eux mêmes, % [sic] eux mêmes -> eux-mêmes, missing hyphen
\[
\begin{tikzcd}[ampersand replacement=\&]
 \& X_1 \arrow[dd] \& X_3 \arrow[l] \arrow[dd] \& \& X_1 \arrow[dd] \& \\
i_0^* :\quad \mathcal{G}\quad\longrightarrow\quad\mathcal{F} : \& \& \& \longmapsto \& \& \\
 \& S \& X_2 \arrow[l] \& \& S \& X_2 \arrow[l]
\end{tikzcd}
\]
\[
\begin{tikzcd}[ampersand replacement=\&]
 \& X_1 \arrow[dd] \& \& \& X_1 \arrow[dd] \& X_1\times_S X_2 \arrow[l] \arrow[dd] \\
i_{0\,*} :\quad \mathcal{F}\quad\longrightarrow\quad\mathcal{G} : \& \& \& \longmapsto \& \& \\
 \& S \& X_2 \arrow[l] \& \& S \& X_2. \arrow[l]
\end{tikzcd}
\]
Ces deux foncteurs sont adjoints l'un de l'autre, et compatibles aux localiseurs $W_{\mathcal{F}}$, $W_{\mathcal{G}}$. Donc les foncteurs déduits de $i_0^*$, $i_{0\,*}$ par passage aux catégories de fractions sont également adjoints (cf. lemme p. 248).

Ainsi, modulo la démonstration des lemmes 1 et 2, on a prouvé ceci :

\printedpage{164}
% head: A. Grothendieck : Dérivateurs, Ch. XII, Édition des Universités de Montpellier II et Paris VII
% foot: Version 30 mars 2003

\msp{page 253}

\textbf{Théorème.} {\itshape Soit $(\mathcal{M},W)$ une catégorie de modèles ($W\subseteq\mathit{Fl}\,\mathcal{M}$) satisfaisant l'axiome de factorisation gauche (p. 240), donnant donc lieu à une bonne théorie des fibres $W$-homotopiques. Soient
\[
\begin{tikzcd}[ampersand replacement=\&, row sep=small, column sep=small]
 \& 1 \& \& \& \& 1 \arrow[dr] \& \\
\Psi=0 \arrow[ur] \arrow[dr] \& \& {} \arrow[r, "i"] \& {} \& \tilde{\Psi}=0 \arrow[ur] \arrow[dr] \& \& 3 \\
 \& 2 \& \& \& \& 2 \arrow[ur] \&
\end{tikzcd}
\]
(immersions ouvertes dans} $\mathrm{Cat}${\itshape ), et considérons le foncteur restriction
\[
i^* : \Ho_W(\tilde{\Psi})\longrightarrow\Ho_W(\Psi).
\]
Ce foncteur admet un adjoint à droite
\[
i_* : \Ho_W(\Psi)\longrightarrow\Ho_W(\tilde{\Psi}).
\]
Ce foncteur $i_*$ se `calcule' ainsi : sa valeur en \begin{tikzcd}[ampersand replacement=\&]
X_1 \arrow[d, "p_1"'] \& \\
S \& X_2 \arrow[l, "p_2"]
\end{tikzcd} est canoniquement isomorphe (dans} $\Ho_W(\tilde{\Psi})${\itshape ) à \begin{tikzcd}[ampersand replacement=\&]
X_1 \arrow[d] \& X_1\times_S X_2 \arrow[l] \arrow[d] \\
S \& X_2 \arrow[l]
\end{tikzcd}, si on suppose que $p_1$ ou $p_2$ est une $W$-fibration. Dans le cas général il se calcule en factorisant $p_1$ en
\[
X_1\xrightarrow{i_1}\bar{X}_1\xrightarrow{p'_1}S,\qquad i_1\in W,\ p'_1\in\mathit{Fib}_W,
\]
et en prenant \begin{tikzcd}[ampersand replacement=\&]
\bar{X}_1 \arrow[d] \& \bar{X}_1\times_S X_2 \arrow[l] \arrow[d] \\
S \& X_2 \arrow[l]
\end{tikzcd}, ou en

\msp{page 254}

procédant de façon symétrique, ou en factorisant simultanement $p_1$ et $p_2$, et prenant % [sic] simultanement -> simultanément, missing accent
\[
\begin{tikzcd}[ampersand replacement=\&]
\bar{X}_1 \arrow[d] \& \bar{X}_1\times_S\bar{X}_2 \arrow[l] \arrow[d] \\
S \& \bar{X}_2. \arrow[l]
\end{tikzcd}
\]}

\dgrblock{dgrb590}{The missing corner can be an entire group}
\begin{dcomment}{examples}{p254:procédant de façon symétrique}{The missing corner can be an entire group}
A groupoid calculation makes the difference between a strict pullback and its homotopy replacement visible without any topological machinery. Work in groupoids with equivalences as $W$, and let $G$ be a group. Write $BG$ for the one-object groupoid with automorphism group $G$. Consider the cospan
\[
 *\longrightarrow BG\longleftarrow *.
\]
Its strict pullback is the terminal groupoid: there is one pair of objects, and the only pair of morphisms is the pair of identities. The homotopy pullback will instead be the discrete groupoid on the set $G$.

To compute it, let $EG$ be the groupoid whose objects are the elements of $G$ and with exactly one morphism from each object to every other object. Define $p:EG\to BG$ by sending the unique morphism $g\to h$ to $hg^{-1}$. The composition formula
\[
 (kh^{-1})(hg^{-1})=kg^{-1}
\]
shows that this is a functor. The inclusion $*\to EG$ selecting the identity element is fully faithful and essentially surjective, hence is an equivalence; its composite with $p$ is the specified map $*\to BG$.

The functor $p$ is an isofibration. Given an object $g$ and an arrow $a$ in $BG$, the arrow $g\to ag$ lies above $a$. Isofibrations of groupoids have the required $W$-fibration property. For completeness, pull back an equivalence $u:C\to D$ over a base along an isofibration $E\to S$. The induced functor on pullbacks is fully faithful because $u$ is. Given an object $(d,e)$ of the target pullback, choose an isomorphism $u(c)\cong d$; its image in $S$ lifts through the isofibration, and the lift changes $e$ to an object lying strictly over the image of $c$. This produces an object in the source pullback isomorphic to $(d,e)$, proving essential surjectivity. The same reasoning applies after any base change.

We can therefore replace the left leg by $*\xrightarrow{\sim}EG\xrightarrow p BG$. In the strict pullback $EG\times_{BG}*$, an object is an element $g\in G$. A morphism from $g$ to $h$ exists precisely when its image $hg^{-1}$ in $BG$ is the identity, which is equivalent to $g=h$. Thus
\[
 EG\times_{BG}*\cong G_{\mathrm{disc}}.
\]
For $G=\mathbb Z/2$, the strict pullback has one connected component while the homotopy pullback has two. For a general group, its elements record the possible isomorphisms between the two images of the chosen object of $*$. The derived completion preserves those comparisons.

This is the cospan version of the groupoid computations already used in the commentary on Chapter~VI, manuscript p.~107. Its purpose here is narrower: it exhibits the actual fourth vertex produced by the theorem and explains why the maps $*\to BG$ must be replaced. The whole completed square, including its two projections and its map to the original cospan after localization, is the value of $i_*$.
\end{dcomment}


% ---- from Der-ChXIIc-30-03-03_p051-055.tex
\printedpage{165}
% head: A. Grothendieck : Dérivateurs, Ch. XII, Édition des Universités de Montpellier II et Paris VII
% foot: Version 30 mars 2003

\textit{La loi fonctorielle s'obtient à l'aide du lemme 1 p. 249, qui permet de se borner aux arguments $\begin{tikzcd}[ampersand replacement=\&] X_1 \arrow[d, "p_1"'] \& \\ S \& X_2 \arrow[l, "p_2"] \end{tikzcd}$ tels que $p_1,p_2\in \mathit{Fib}_W$, de sorte que l'on trouve un foncteur}
\[
\begin{tikzcd}[ampersand replacement=\&]
X_1 \arrow[d] \& \\
S \& X_2 \arrow[l]
\end{tikzcd}
\qquad\longmapsto\qquad
\begin{tikzcd}[ampersand replacement=\&]
X_1 \arrow[d] \& X_1\times_S X_2 \arrow[l] \arrow[d] \\
S \& X_2, \arrow[l]
\end{tikzcd}
\]
\textit{compatible aux localiseurs, et passant aux catégories de fractions} $\Ho_W(\Psi)$ \textit{et} $\Ho_W(\tilde{\Psi})$ \textit{respectivement.}

\textbf{NB} En procédant de façon similaire, plus simple (sans besoin d'un recours au lemme 1), on trouve l'existence d'un \textit{adjoint à gauche} $i_!$ de $i^* : \Ho_W(\Psi)\longrightarrow \Ho_W(\tilde{\Psi})$, sous réserve de l'existence d'un objet initial $\varphi$ dans $\mathcal{M}$. Cet foncteur $i_!$ est donné par % [sic] Cet foncteur -> Ce foncteur, "foncteur" begins with a consonant

\msp{page 255}
\[
i_!\left(
\begin{tikzcd}[ampersand replacement=\&]
X_1 \arrow[d] \& \\
S \& X_2 \arrow[l]
\end{tikzcd}
\right)
\qquad = \qquad
\left(
\begin{tikzcd}[ampersand replacement=\&]
X_1 \arrow[d] \& \varphi \arrow[l] \arrow[d] \\
S \& X_2 \arrow[l]
\end{tikzcd}
\right),
\]
la loi fonctorielle étant évidente.

\dgrblock{dgrb591}{The adjunction determines the completion}
\begin{dcomment}{exposition}{p254}{The adjunction determines the completion}
The adjunction is what makes Grothendieck's construction a canonical completion of a cospan. On the good-diagram categories, restriction $r:\mathcal G\to\mathcal F$ has right adjoint $c$, where $c$ supplies the strict pullback. A morphism from a square $Q$ to $c(d)$ is uniquely determined by its three components on the cospan $r(Q)$: the fourth component exists uniquely by the pullback universal property. This gives the adjunction bijection before localization.

The counit $rc\to1_{\mathcal F}$ is the identity on the original three vertices. The unit at a square with fourth vertex $T$ is the map to its completion, whose fourth component is
\[
 T\longrightarrow A\times_S B.
\]
At a completed square this map is the identity, and restricting the unit to its three original vertices also gives the identity. These observations verify both triangle identities directly. Restriction preserves objectwise weak equivalences; the preceding base-change calculation proves that completion does too. The localization lemma on ms.~248 therefore descends the adjunction. Once Lemmas~1 and~2 identify the localized good-diagram categories with the required diagram categories, it is exactly $i^*\dashv i_*$, with invertible counit $i^*i_*\to1$.

The formula using only one fibrant leg follows by an explicit comparison, not by a symmetry convention. If $A\to S$ is already a $W$-fibration, replace only $B$ by $\bar B$. The map
\[
 A\times_S B\longrightarrow A\times_S\bar B
\]
is in $W$, as in the block at ms.~242, and forms an objectwise weak equivalence between the completed square and one with two good legs. The latter computes $i_*$. For an arbitrary cospan, first replace $A$; the same argument compares the one-leg and two-leg formulas. Interchanging $A$ and $B$ gives the other formula. These comparisons concern objects of $\Ho_W(\tilde\Psi)$ and retain all the projection maps.

There is also the left adjoint mentioned in the concluding NB. If $\mathcal M$ has an initial object $\varphi$, extending a cospan by putting $\varphi$ at vertex~$3$ is left adjoint to restriction already on ordinary diagram categories. The fourth component of any extension of a map of cospans is the unique map $\varphi\to T$. This extension preserves objectwise weak equivalences, since its fourth component on every morphism is $1_\varphi$, so the adjunction descends by the same lemma. It needs no fibrant replacement. The displayed type of $i^*$ in that NB is reversed in the edition; its correct type remains $\Ho_W(\tilde\Psi)\to\Ho_W(\Psi)$. Likewise the $\Phi,\tilde\Phi$ used on ms.~251--252 denote the $\Psi,\tilde\Psi$ of the theorem, not a second pair of diagram shapes.
\end{dcomment}


\dgrblock{dgrb592}{Two extensions of the same cospan}
\begin{dcomment}{exposition}{p255:la loi fonctorielle étant évidente.}{Two extensions of the same cospan}
Grothendieck has now reached the point at which a homotopy fibre can be characterized by its relation to all diagrams, rather than by a particular construction of paths. The two extensions of the three-vertex diagram make this distinction visible. Restriction has the direction
\[
 i^*: \Ho_W(\tilde\Psi)\longrightarrow\Ho_W(\Psi).
\]
The reversed direction in the preceding NB is an edition slip. Its left adjoint, discussed in the preceding block, fills the fourth vertex with an initial object; its right adjoint fills that vertex with a homotopy pullback. These operations answer different universal problems on the same cospan.

For the right adjoint, the fourth vertex must instead receive a compatible pair of maps. Once the two legs to $S$ are $W$-fibrations, that universal object is $X_1\times_S X_2$. The extra work in Lemmas 1 and 2 is what allows this calculation on suitably replaced diagrams to determine a functor and an adjunction on the original localization. The homotopy limit of the cospan is then obtained by evaluating this completed square at vertex 3.

The indexing convention explains the apparently reversed endpoint in the following lemma. An object of $\mathcal M(I)$ is a functor $I^{\mathrm{op}}\to\mathcal M$. If $t$ is terminal in $I$, each arrow $i\to t$ induces $F(t)\to F(i)$. A natural transformation from the constant diagram on $A$ to $F$ is consequently determined by its component $A\to F(t)$, and every such component gives a natural transformation. Thus
\[
 \operatorname{Hom}_{\mathcal M(I)}(p^*A,F)
 \cong \operatorname{Hom}_{\mathcal M}(A,F(t)).
\]
Evaluation at $t$ is a right adjoint. The displayed formula survives localization because both functors preserve the designated weak equivalences. In particular, the object constructed here is universal for maps of coherent cospans. There is no assertion that taking an ordinary pullback in $\mathcal M[W^{-1}]$ would recover it.
\end{dcomment}


Ce n'est pas encore la fin du sorite! Il faut encore interpréter le foncteur composé
\[
\Ho_W(\Psi)\xrightarrow{\ i_*\ }\Ho_W(\tilde{\Psi})\xrightarrow{\ \varepsilon_3^*\ }\Ho_W(e)=\Ho_W,
\]
où
\[
\varepsilon_3 : e\xrightarrow{\ \sim\ }\{3\}
\qquad\longrightarrow\qquad
\begin{tikzcd}[ampersand replacement=\&, column sep=small, row sep=small]
\& 1 \arrow[dl] \& \\
3 \& \& 0 \arrow[ul] \arrow[dl] \\
\& 2 \arrow[ul] \&
\end{tikzcd}
\]
est l'inclusion du sommet 3 dans $\tilde{\Psi}$. Si $\Ho_W(\Psi)$ est interprété, conformément au lemme 1, comme $\mathcal{F}W_{\mathcal{F}}^{-1}$, ce foncteur composé n'est autre que
\[
\begin{tikzcd}[ampersand replacement=\&]
X_1 \arrow[d] \& \\
S \& X_2 \arrow[l]
\end{tikzcd}
\qquad\longmapsto\qquad X_1\times_S X_2
\]
\printedpage{166}
% head: A. Grothendieck : Dérivateurs, Ch. XII, Édition des Universités de Montpellier II et Paris VII
% foot: Version 30 mars 2003
(pour $\left(\begin{tikzcd}[ampersand replacement=\&] X_1 \arrow[d, "p_1"'] \& \\ S \& X_2 \arrow[l, "p_2"] \end{tikzcd}\right)\in \mathrm{Ob}\,\mathcal{F}$, i.e. $p_1,p_2\in\Fib_W$.)

On veut interpréter ce foncteur comme un $\varprojlim$ $W$-homotopique, i.e. vu comme un `produit fibré $W$-homotopique'. De façon précise, on a ceci :

\msp{page 256}
\textbf{Corollaire 1.} \textit{Considérons $p_\Psi : \Psi\longrightarrow e$, et le foncteur}
\[
p_\Psi^* : \Ho_W(e)=\Ho_W\longrightarrow\Ho_W(\Psi),
\]
\textit{déduit par passage aux catégories de fractions du foncteur}
\[
\begin{array}{ccc}
\mathcal{M} & \longrightarrow & \mathcal{M}(\Psi) \\
X & \longrightarrow & \left(
\begin{tikzcd}[ampersand replacement=\&]
X \arrow[d] \& \\
X \& X \arrow[l]
\end{tikzcd}
\right).
\end{array}
\]
\textit{Ce foncteur admet le foncteur $\varepsilon_3^*i_*$ précédent comme adjoint à droite, i.e.}
\[
p_{\Psi\,*}=\varepsilon_3^*i_* : \Ho_W(\Psi)\longrightarrow\Ho_W.
\]

\textsc{Démonstration}. On peut regarder $p_\psi$ comme le composé % [sic] p_\psi -> p_\Psi, the subscript is the capital \Psi everywhere else
\[
\Psi\xrightarrow{\ i\ }\tilde{\Psi}\xrightarrow{\ p_{\tilde{\Psi}}\ }e,
\]
et on a donc
\[
p_\Psi^*\simeq i^*(p_{\tilde{\Psi}})^*.
\]
Donc l'existence et la détermination d'un adjoint à droite pour $p_\Psi^*$ résultera de l'existence et de la détermination d'adjoints à droite $i_*$, $(p_{\tilde{\Psi}})_*$ de $i^*$, $(p_{\tilde{\Psi}})^*$, via
\[
(p_\Psi)_*\simeq(p_{\tilde{\Psi}})_*i_*.
\]

\msp{page 257}
Donc le corollaire 1 résulte de la formule
\[
p_{\tilde{\Psi}\,*}=\varepsilon_3^*,
\]
qui elle-même est un cas particulier du

\printedpage{167}
% head: A. Grothendieck : Dérivateurs, Ch. XII, Édition des Universités de Montpellier II et Paris VII
% foot: Version 30 mars 2003

\textbf{Lemme.} \textit{Soit $I$ dans} Cat\textit{, avec objet final $e$, considérons l'inclusion}
\[
\varepsilon : e\longrightarrow I,\qquad \textit{et } p : I\longrightarrow e.
\]
\textit{Soit $(\mathcal{M},W)$ une catégorie de modèles, et considérons}
\[
\Ho_W(I)\xleftarrow{\ p^*\ }\Ho_W.
\]
\textit{Ce foncteur admet un adjoint à gauche} \texttt{[plutôt à droite]} \textit{$p_*$, lequel est isomorphe à $\varepsilon^*$.}

\textsc{Démonstration}. Considérons le couple de foncteurs adjoints
\[
\begin{tikzcd}[ampersand replacement=\&, column sep=large]
\mathcal{M}(I) \arrow[r, shift right, "p_*=\varprojlim\limits_I"'] \& \mathcal{M} \arrow[l, shift right, "p^*"']
\end{tikzcd}\ ;
\]
comme $I$ admet l'objet final $e$, le foncteur $p_*$ est isomorphe à $\varepsilon^*$. Ces foncteurs $p^*$, $\varepsilon^*$ sont compatibles aux localiseurs $W(I)$, $W$, donc on conclut par le lemme page 248.

Il ne reste plus qu'à prouver les lemmes 1 et 2. Je le ferai demain.

\msp{page 258}
Aujourd'hui, j'ai envie encore de donner les sorites des carrés $W$-cartésiens dans une

\textbf{Proposition 19 :} \textit{Soit $\begin{tikzcd}[ampersand replacement=\&] X_1 \arrow[d, "p_1"'] \& X_3 \arrow[l, "f_1"'] \arrow[d, "f_2"] \\ S \& X_2 \arrow[l, "p_2"] \end{tikzcd}$ un carré commutatif dans $\mathcal{M}$. Considérons des factorisations de $p_1$, $p_2$ en}
\[
X_1\xrightarrow{\ i_1\ }\bar{X}_1\xrightarrow{\ \bar{p}_1\ }S,\qquad X_2\xrightarrow{\ i_2\ }\bar{X}_2\xrightarrow{\ \bar{p}_2\ }S,
\]
\textit{$i_1,i_2\in W$, $\bar{p}_1,\bar{p}_2\in F_W$} \texttt{[$=$ Fib$_W$]}\textit{, et considérons les flèches induites par $X_3\longrightarrow X_1\times_S X_2$ (déduite de $Q$) via $\begin{tikzcd}[ampersand replacement=\&, row sep=small] \& \bar{X}_1\times_S X_2 \\ X_1\times_S X_2 \arrow[ur] \arrow[r] \arrow[dr] \& \bar{X}_1\times_S\bar{X}_2 \\ \& X_1\times_S\bar{X}_2 \end{tikzcd}$, de sorte qu'on trouve un diagramme commutatif}
\[
\begin{tikzcd}[ampersand replacement=\&, row sep=small]
\& \bar{X}_1\times_S X_2 \arrow[d, "j_2"'] \\
X_3 \arrow[ur, "\pi_1"] \arrow[r, "\pi_0"] \arrow[dr, "\pi_2"] \& \bar{X}_1\times_S\bar{X}_2 \\
\& X_1\times_S\bar{X}_2, \arrow[u, "j_1"]
\end{tikzcd}
\]
\textit{où $j_1=i_1\times_S\mathrm{id}_{\bar{X}_2}$ et $j_2=\mathrm{id}_{\bar{X}_1}\times i_2$ sont dans $W$ (car $i_1$ et $i_2$ le sont, et $\bar{X}_1$ et $\bar{X}_2$ sont dans $\underline{\mathrm{Fib}}_W S$). On a donc}
\[
(\pi_1\in W)\qquad\iff\qquad(\pi_2\in W)\qquad\qquad(\iff\ \pi_0\in W),
\]
\textit{ce qui prouve que la condition $\pi_1\in W$, qui semble dépendre du choix de la factorisation par $i_1$, n'en dépend pas, et idem pour $\pi_2\in W$, et pour $\pi_0\in W$ (qui semble dépendre à la fois de $i_1$, et de $i_2$).}

\dgrblock{dgrb593}{Independence of replacements and invariance of the test}
\begin{dcomment}{exposition}{p258:ce qui prouve que la condition}{Independence of replacements and invariance of the test}
Grothendieck wants the property of a square to be independent of the auxiliary factorization used to recognize it. Proposition 19 achieves this with a particularly economical comparison: one changes the first leg while holding the second fixed, then exchanges their roles. We spell out the argument because it also supplies the invariance used in the next pasting proof.

Write a square as $D\to A$, $D\to B$ over the cospan $A\to S\leftarrow B$. Choose $A\xrightarrow{a}\bar A\to S$ and $B\xrightarrow{b}\bar B\to S$ as in the proposition. The two arrows
\[
 \bar A\times_S B\longrightarrow\bar A\times_S\bar B,
 \qquad
 A\times_S\bar B\longrightarrow\bar A\times_S\bar B
\]
are in $W$, since each pulls back a member of $W$ along a $W$-fibration. Thus the three comparison maps from $D$ are simultaneously in $W$. Now choose a second factorization of $A\to S$, retaining the chosen factorization of $B\to S$. Both first-leg tests are equivalent to the very same test $D\to A\times_S\bar B$. They therefore agree. The argument with $A$ and $B$ interchanged proves independence of the other choice. No map between the two factorizations of $A$ has been assumed.

For the pasting argument we also need the following verification. Suppose a map of squares $Q\to Q'$ is in $W$ at all four vertices. Denote the cospans by $A\to S\leftarrow B$ and $A'\to S'\leftarrow B'$, and choose a factorization
\[
 A'\xrightarrow{a'}\widehat A'\xrightarrow{q'}S',
 \qquad a'\in W,\quad q'\in\Fib_W.
\]
Set $T=S\times_{S'}\widehat A'$. The projection $T\to\widehat A'$ is in $W$, being the pullback of $S\to S'$ along $q'$. The induced $A\to T$ is in $W$ by two-out-of-three, since its composite to $\widehat A'$ is $A\to A'\xrightarrow{a'}\widehat A'$. Moreover $T\to S$ is a $W$-fibration. This is an admissible replacement of the first leg of $Q$.

With these choices the comparison of the two pullback objects is
\[
 T\times_S B\cong\widehat A'\times_{S'}B
 \longrightarrow\widehat A'\times_{S'}B'.
\]
It is in $W$, since it is obtained by pulling back $B\to B'$ along the pullback of $q'$. The square relating these objects to $D\to D'$ commutes. Its vertical arrows are in $W$, so two-out-of-three shows that the comparison for $Q$ is in $W$ if and only if that for $Q'$ is. This proves invariance under objectwise weak equivalences, and hence under zigzags of them, using only the factorization hypothesis and the intrinsic definition of $\Fib_W$. Strong saturation, functorial choices of factorization, and the deferred adjunction lemmas are unnecessary for this particular assertion.
\end{dcomment}


\printedpage{168}
% head: A. Grothendieck : Dérivateurs, Ch. XII, Édition des Universités de Montpellier II et Paris VII
% foot: Version 30 mars 2003

\msp{page 259}
\textbf{Définition.} \textit{Quand ces conditions sont satisfaites, on dit que le carré $Q$ est $W$-cartésien.}

\textbf{Corollaire 1.} \textit{Si $Q$ est $W$-cartésien, alors l'image de $Q$ dans} $\Ho_W(\tilde{\Psi})$ \textit{est dans l'image essentielle de}
\[
i_* : \Ho_W(\Psi)\longrightarrow\Ho_W(\tilde{\Psi}),
\]
\textit{i.e. la flèche canonique dans} $\Ho_W(\tilde{\Psi})$
\[
Q\longrightarrow i_*i^*Q
\]
\textit{est} \texttt{[un]} \textit{isomorphisme. La réciproque est vraie si $W$ est fortement saturé.}

\dgrblock{dgrb594}{The unit and the role of strong saturation}
\begin{dcomment}{exposition}{p259:La réciproque est vraie si}{The unit and the role of strong saturation}
Grothendieck now compares his concrete test with the intrinsic language of homotopy Kan extension. This comparison separates two tasks which coincide in the familiar examples: constructing an isomorphism after localization, and recognizing that a specified original arrow belongs to $W$.

Let $Q^r$ be obtained by replacing $X_1\to S$ and $X_2\to S$ by $\bar X_1\to S$ and $\bar X_2\to S$, while keeping $X_3$ and $S$. There is an objectwise weak equivalence $Q\to Q^r$. The unit on $Q^r$ is represented by the map to the completed pullback square which is the identity on $\bar X_1$, $\bar X_2$, and $S$, and is
\[
 \pi_0:X_3\longrightarrow\bar X_1\times_S\bar X_2
\]
at the fourth vertex. If $\pi_0\in W$, this map of squares becomes an isomorphism. Naturality of the unit then gives the assertion for $Q$.

Conversely, if the unit on $Q$ is invertible, the unit on $Q^r$ is invertible. Evaluation at vertex 3 shows that $\pi_0$ is inverted by $\mathcal M\to\mathcal M[W^{-1}]$. Strong saturation says exactly that this forces $\pi_0\in W$. Two-out-of-three by itself makes no such assertion about every map inverted by localization. This is the reason for the qualification in Corollary 1, rather than an incidental technical precaution.

Since $i^*i_*\cong\mathrm{id}$, the right adjoint is fully faithful, and the usual adjunction argument identifies its essential image with the squares whose unit is invertible. In modern derivator terminology this essential-image condition defines cartesian squares; the concrete $W$ test agrees with it under the saturation hypothesis just explained. See \citet[Proposition 1.20 and Lemma 1.21]{Groth2013} for the fully faithful Kan-extension formulation. That formulation presupposes the construction of the adjunction, so it does not by itself discharge the two deferred localization lemmas in this manuscript.
\end{dcomment}


\textbf{Corollaire 2.} \textit{Considérons un diagramme commutatif dans $\mathcal{M}$}
\[
\begin{tikzcd}[ampersand replacement=\&]
X_0' \arrow[d] \& X_1' \arrow[l] \arrow[d] \arrow[dl, phantom, "\textit{1}"] \& X_2' \arrow[l] \arrow[d] \arrow[dl, phantom, "\textit{2}"] \\
X_0 \& X_1 \arrow[l] \& X_2, \arrow[l]
\end{tikzcd}
\tag*{(D)}
\]
\textit{d'où deux carrés commutatifs 1, 2, et un carré composé}
\[
\begin{tikzcd}[ampersand replacement=\&]
X_0' \arrow[d] \& X_2' \arrow[l] \arrow[d] \arrow[dl, phantom, "\textit{3}"] \\
X_0 \& X_2. \arrow[l]
\end{tikzcd}
\]
\textit{Supposons le carré 1 $W$-cartésien. Alors le carré 2 est $W$-cartésien si et seulement si le carré 3 l'est.}

\textsc{Démonstration}. Considérons une factorisation de $\begin{tikzcd}[ampersand replacement=\&, row sep=small] X_0' \arrow[d] \\ X_0 \end{tikzcd}$ en
\[
\begin{tikzcd}[ampersand replacement=\&, column sep=small, row sep=small]
\bar{X}_0' \arrow[dr] \& X_0' \arrow[l] \arrow[d] \\
\& X_0
\end{tikzcd}
\]
avec $(X_0'\longrightarrow\bar{X}_0')\in W$, et $(\bar{X}_0'\longrightarrow X_0)\in\Fib_W$. On peut alors remplacer dans (D) $X_0'$ par $\bar{X}_0'$, obtenant un diagramme $(\bar{D})$, et on voit aussitôt
\printedpage{169}
% head: A. Grothendieck : Dérivateurs, Ch. XII, Édition des Universités de Montpellier II et Paris VII
% foot: Version 30 mars 2003
\msp{page 260}
que chacun des carrés 1(D), 2(D), 3(D) est $W$-cartésien si et seulement si le carré correspondant $1(\bar{D})$, $2(\bar{D})$ ou $3(\bar{D})$ l'est. Cela montre qu'on peut supposer $X_0\longrightarrow S$ dans $\Fib_W S$. Considérons alors le diagramme commutatif déduit de (D)
\[
\begin{tikzcd}[ampersand replacement=\&, column sep=small]
X_0' \arrow[d] \& \tilde{X}_1'\overset{\text{déf}}{=}X_0'\times_{X_0}X_1 \arrow[l] \arrow[d] \& X_1' \arrow[l, "\alpha"'] \& X_2' \arrow[l] \arrow[d] \\
X_0 \& X_1 \arrow[l] \& \& X_2, \arrow[ll]
\end{tikzcd}
\]
où $\alpha\in W$ en vertu de l'hypothèse $W$-cartésienne sur le carré 1. On voit encore que le carré 2 est $W$-cartésien si et seulement si le carré correspondant avec $X_1'$ remplacé par $\tilde{X}_1'$ l'est, quant aux carrés 1, 3, ils ne changent pas par cette modification. Cela montre qu'on peut supposer le carré 1 \textit{cartésien.} Considérons alors le diagramme
\[
\begin{tikzcd}[ampersand replacement=\&]
X_0' \arrow[d, "\mathrm{Fib}_W"'] \& X_1' \arrow[l] \arrow[d, "\mathrm{Fib}_W"] \arrow[dl, phantom, "\text{cart.}"] \& \tilde{X}_2'\overset{\text{déf}}{=}X_1'\times_{X_1}X_2 \arrow[l] \arrow[d, "\mathrm{Fib}_W"] \arrow[dl, phantom, "\text{cart.}"] \& X_2' \arrow[l, "\beta"'] \\
X_0 \& X_1 \arrow[l] \& X_2. \arrow[l]
\end{tikzcd}
\]
Alors la proposition 19 et la définition montrent que le carré 2 est $W$-cartésien si et seulenment si $\beta\in W$, et de même pour le carré 3. Ce qui achève la démonstration. % [sic] seulenment -> seulement, misspelling

\dgrblock{dgrb595}{Pasting reduces two questions to the same arrow}
\begin{dcomment}{exposition}{p260:Ce qui achève la démonstration.}{Pasting reduces two questions to the same arrow}
Grothendieck's proof makes pasting effective by arranging that the first square is an actual pullback along a $W$-fibration. The apparently immediate replacements in the proof are justified by the objectwise invariance established in our preceding block. The printed phrase $X_0\to S\in\Fib_W S$ should be read as $X'_0\to X_0\in\Fib_W$; it is the latter arrow that has just been factored.

Write $E\to X_0$ for a fibrant replacement of $X'_0\to X_0$. Replacing $X'_0$ by $E$ changes neither relevant cartesianness test. Because the first square is $W$-cartesian, its comparison
\[
 X'_1\longrightarrow E\times_{X_0}X_1
\]
is in $W$. We may therefore replace $X'_1$ by $E_1=E\times_{X_0}X_1$, preserving the test for the second square. The first square is now strictly cartesian, and $E_1\to X_1$ is a $W$-fibration.

Both remaining conditions are consequently tested by the one arrow
\[
 X'_2\longrightarrow E_1\times_{X_1}X_2
       \cong E\times_{X_0}X_2.
\]
Using the expression on the left tests the second square; using the expression on the right tests the outer rectangle. The two conditions are identical. This proves both the composition implication and its indicated cancellation, without assuming that arbitrary base change preserves $W$.

The result became a basic part of the calculus of derivators \citep[Proposition 3.13]{Groth2013}. Its orientation matters: in the manuscript the horizontal arrows point towards the left, so its left-hand square is the square adjacent to the common target. One cannot replace the hypothesis by either small square indiscriminately in an unstable theory.
\end{dcomment}


\textbf{Corollaire 2.} \textit{Soit $\begin{tikzcd}[ampersand replacement=\&] X_1 \arrow[d, "p_1"'] \& X_3 \arrow[l, "f_2"'] \arrow[d, "f_1"] \\ S \& X_2 \arrow[l, "p_2"] \end{tikzcd}$ un carré $W$-cartésien. Alors}
\[
\begin{aligned}
p_1\in W\quad&\Longrightarrow\quad f_1\in W\\
p_2\in W\quad&\Longrightarrow\quad f_2\in W,
\end{aligned}
\]
\textit{plus précisément, si $p_i\in W$, alors $Q$ est $W$-cartésien si et seulement si $f_i\in W$.}

\msp{page 261}
Supposons $p_1\in W$. Alors dans la factorisation de $p_1$ en $X_1\longrightarrow\bar{X}_1\longrightarrow S$, on peut prendre $\bar{X}_1=S$. Donc le carré est $W$-cartésien si et seulement si
\[
(X_3\longrightarrow\bar{X}_1\times_S X_2=X_2)\in W,
\]
i.e. $f_1\in W$, q.e.d.

Bien sûr, on ne peut en général affirmer les implications inverses : $f_i\in W$ n'implique nullement $p_i\in W$.

\dgrblock{dgrb596}{A nonempty counterexample to the converse}
\begin{dcomment}{examples}{p261:Bien sûr, on ne peut en général affirmer}{A nonempty counterexample to the converse}
Grothendieck's final caution about cartesian squares concerns a loss of information under restriction to part of the base. A finite example already shows the issue, without empty objects or any higher homotopy calculation.

Take $\mathcal M=\mathrm{Set}$ and let $W$ be the bijections. Every map is an intrinsic $W$-fibration, and the identity followed by the given map provides the required factorization. Put
\[
 S=\{0,1\},\quad X_1=\{a,b,c\},\quad X_2=\{0\},
 \qquad p_1(a)=0,\quad p_1(b)=p_1(c)=1,
\]
and let $p_2$ be the inclusion of $\{0\}$ in $S$. The pullback is $X_3=\{a\}$. Its projection $f_1:X_3\to X_2$ is a bijection, whereas $p_1:X_1\to S$ is not. The square is strictly cartesian and hence $W$-cartesian.

Only the fibre over 0 has been tested. Nothing in the upper projection records the two elements over 1. The forward assertion in the source has no such problem: if $p_1\in W$, use the factorization $X_1\xrightarrow{p_1}S\xrightarrow{\mathrm{id}}S$, and the cartesianness comparison is precisely $f_1$. The later appendix will ask for a converse to the analogous pushout statement. That converse is an additional conservativity condition, rather than a formal consequence of the existence of homotopy pushouts.
\end{dcomment}


% ---- from Der-ChXIIc-30-03-03_p056-060.tex
\printedpage{170}
% head: A. Grothendieck : Dérivateurs, Ch. XII, Édition des Universités de Montpellier II et Paris VII

Tous ces résultats, et la notion de carré $W$-cartésien, se dualisent (moyennant l'axiome idoine de factorisation à droite, cf. page 240) en une théorie des sommes amalgamées $W$-homotopiques, et des carrés $W$-cocartésiens. Cette fois on travaille avec la catégorie
\[
\Phi \;=\; \left(
\begin{tikzcd}[ampersand replacement=\&, row sep=small, column sep=small]
 \& 1 \arrow[dl] \\
0 \& \\
 \& 2 \arrow[ul]
\end{tikzcd}
\right)
\qquad\quad \overset{j}{\hookrightarrow} \qquad\quad
\tilde{\Phi} \;=\; \left(
\begin{tikzcd}[ampersand replacement=\&, row sep=small, column sep=small]
 \& 1 \arrow[dl] \& \\
0 \& \& 3 \arrow[ul] \arrow[dl] \\
 \& 2 \arrow[ul] \&
\end{tikzcd}
\right),
\]
$j$ étant une immersion fermée. On a un foncteur restriction
\[
j^* : \Ho_W(\tilde{\Phi}) \longrightarrow \Ho_W(\Phi),
\]
dont on construit un adjoint à gauche
\[
j_! : \Ho_W(\Phi) \longrightarrow \Ho_W(\tilde{\Phi}).
\]
Et on trouve une notion naturelle de
\msp{page 262}
carré $W$-cocartésien, qui paraît irréprochable à tout points de vue, et d'une simplicité parfaite. % [sic] tout points -> tous points, agreement of "tout"

Cela amène à revoir la notion en question dans Cat, où elle avait l'air très astucieuse, profonde et un peu obscure, avec des intégrales et tout ça. C'est le moment de se rendre compte que `ce n'est que ça'. Donc de prouver la

\textbf{Proposition 20.} \textit{Supposons} $\mathcal{M} = \mathrm{Cat}$\textit{, $W$ un localiseur fondamental satisfaisant W(7 bis). Considérons un carré commutatif dans} Cat
\[
Q = \left(
\begin{tikzcd}[ampersand replacement=\&]
S \arrow[r] \arrow[d] \& X_1 \arrow[d] \\
X_2 \arrow[r] \& X_3
\end{tikzcd}
\right).
\]
\textit{Pour qu'il soit $W$-cartésien} \texttt{[plutôt} $W$\texttt{-cocartésien]} \textit{au sens de la théorie développée ici (pour des catégories de modèles générales), il faut et il suffit que}
\[
\varphi_Q : \quad \mathcal{X}_Q = \left(
\begin{tikzcd}[ampersand replacement=\&, row sep=small, column sep=small]
 \& X_1 \arrow[dl] \\
\int S \& \\
 \& X_2 \arrow[ul]
\end{tikzcd}
\right) \qquad\quad \longrightarrow \qquad\quad X_3
\]

\texttt{[plutôt}
\[
\varphi_Q : \quad \mathcal{X}_Q = \left(
\begin{tikzcd}[ampersand replacement=\&, row sep=small, column sep=small]
 \& X_1 \\
\int S \arrow[ur] \arrow[dr] \& \\
 \& X_2
\end{tikzcd}
\right) \qquad\quad \longrightarrow \qquad\quad X_3\text{\texttt{]}}
\]
\printedpage{171}
% head: A. Grothendieck : Dérivateurs, Ch. XII, Édition des Universités de Montpellier II et Paris VII
\textit{soit dans $W$.} (NB L'intérêt de ce critère est qu'il est `intrinsèque' à $Q$, ne faisant pas appel à une factorisation de $\beta_1$ ou de $\beta_2$ \texttt{[où} $\beta_i : S \longrightarrow X_i$\texttt{]} en $S \xrightarrow{\;\in \Cof_W\;} \tilde{X}_i \xrightarrow{\in W} X_i$. Cf. application intéressante dans section V 3.4., p. 146 ff., th. 4 et corollaires.)

\dgrblock{dgrb597}{The integral as a categorical homotopy pushout}
\begin{dcomment}{later}{p262:L'intérêt de ce critère est qu'il est}{The integral as a categorical homotopy pushout}
Grothendieck wants to identify the abstract operation just constructed with the categorical integral used earlier in the chapter. Proposition 20 says that the older construction was already providing a model for the same homotopy pushout. Its advantage is that the comparison can be read directly from the original square.

Here $\Phi$ has arrows $1\to0\leftarrow2$, and the diagram of categories is contravariant: its data are $X_1\leftarrow S\to X_2$. The associated fibred Grothendieck construction has objects $(i,x)$, with $x$ an object of the fibre at $i$. For $\alpha:i\to j$, a morphism over $\alpha$ is a map $x\to F(\alpha)(y)$ in the fibre at $i$. A compatible cocone to $X_3$ gives the displayed functor $\varphi_Q$.

For ordinary weak homotopy equivalences, this interpretation is the content of Thomason's homotopy-colimit theorem: the classifying space of the integral models the homotopy colimit of the classifying spaces of the diagram \citep[Theorem 1.2]{Thomason1979}. We checked its statement and a later proof in \citet[Theorem 26.8 and Corollary 26.9]{ChacholskiScherer2001}; the verification note distinguishes this from inspection of the original 1979 proof. In the present shape it gives
\[
 \left|N\mathcal X_Q\right|
 \simeq |NX_1|\amalg^{h}_{|NS|}|NX_2|.
\]
The fibred version follows from the usual covariant theorem by applying it to the opposite fibre categories on $\Phi^{\mathrm{op}}$ and then taking the opposite total category. Taking opposites leaves classifying spaces canonically homeomorphic. The equivalence explains the notation, but the manuscript's argument addresses a general $W$ under its stated axioms, rather than only $W_\infty$.

Its proof has one directional slip worth resolving before following the reduction. The transformation $\alpha:\tilde Q\to Q$ induces
\[
 \lambda:\mathcal X_{\tilde Q}\longrightarrow\mathcal X_Q,
 \qquad \varphi_{\tilde Q}=\varphi_Q\lambda.
\]
The edition reverses the horizontal arrow in this triangle. With the correct direction, the fibrewise weak-equivalence assertion gives $\lambda\in W$, and two-out-of-three gives the required equivalence of tests. Footnote 134 expressly identifies the strong localization axiom used for this step; it should not be silently dropped when $W$ is left arbitrary.

After both legs have been replaced by Dwyer immersions, put $P=X_1\amalg_SX_2$. The earlier result on pushouts of these immersions supplies $\mathcal X_Q\to P$ in $W$. The map to $X_3$ factors as $\mathcal X_Q\to P\xrightarrow{\pi}X_3$, so $\varphi_Q\in W$ if and only if $\pi\in W$. The integral and the abstract replacement criterion therefore agree. None of this requires the reverse assertion that a homotopy pushout reflects weak equivalences in an attached leg; that separate condition enters the appendix.
\end{dcomment}


\msp{page 263}
\textsc{Démonstration.} Considérons une factorisation de $S \longrightarrow X_1$ en
\[
S \xrightarrow{i_1} \tilde{X}_1 \xrightarrow{p_1} X, % [sic] X -> X_1, the factorisation is of S -> X_1
\]
avec $i_1$ immersion ouverte de Dwyer, $p_1 \in W$. Donc on a un carré commutatif
\[
\tilde{Q} = \left(
\begin{tikzcd}[ampersand replacement=\&]
S \arrow[r, hook, "\mathrm{Dw}"] \arrow[d] \& \tilde{X}_1 \arrow[d] \\
X_2 \arrow[r] \& X_3
\end{tikzcd}
\right)
\]
et un homomorphisme $\tilde{Q} \xrightarrow{\alpha} Q$ qui est l'identité sur les trois sommets distincts de $\tilde{X}_1$, et qui est $p_1$ en $\tilde{X}_1$. On en déduit un diagramme commutatif
\[
\begin{tikzcd}[ampersand replacement=\&, row sep=small, column sep=small]
\mathcal{X}_Q \arrow[rr, "\mathcal{X}_\alpha = \lambda"] \arrow[dr, "\varphi_Q"'] \& \& \mathcal{X}_{\tilde{Q}} \arrow[dl, "\varphi_{\tilde{Q}}"] \\
 \& X_3, \&
\end{tikzcd}
\]
et on a $\lambda \in W$ (car c'est une $\Phi$-flèche de catégories fibrées sur $\Phi$, qui induit une $W$-équivalence sur chacune de trois fibres.) Donc $\varphi_Q \in W \iff \varphi_{\tilde{Q}} \in W$, d'autre part il est clair que $Q$ est $W$-cocartésien si et seulement si $\tilde{Q}$ l'est. Donc on peut supposer que $S \hookrightarrow X_1$ est \texttt{[une]} immersion ouverte (de Dwyer) (\footnote[134]{On utilise ici W(3) (forme forte), i.e. Loc(3 bis) avec les notations de XVI. Mais on peut s'en passer, cf. plus bas.}). Appliquant ceci
\msp{page 264}
en changeant les rôles de $X_1$ et $X_2$, on voit qu'on peut supposer de plus que $S \longrightarrow X_2$ est également \texttt{[une]} immersion ouverte de Dwyer. Considérons alors
\[
\begin{tikzcd}[ampersand replacement=\&]
S \arrow[r, hook] \arrow[d, hook] \arrow[dr, phantom, "Q'"] \& X_1 \arrow[d, hook] \& \\
X_2 \arrow[r, hook] \& X_1 \amalg_S X_2 \arrow[r, "\pi"] \& X_3 \\
 \& \mathcal{X}_Q \arrow[u, "\varphi_{Q'}"] \arrow[ur, "\varphi_Q"'] \&
\end{tikzcd}
\]
\printedpage{172}
% head: A. Grothendieck : Dérivateurs, Ch. XII, Édition des Universités de Montpellier II et Paris VII
On a $\varphi_Q = \pi\varphi_{Q'}$, d'autre part on sait que $\varphi_{Q'} \in W$ (cor. 3, p. 149). Donc
\[
\pi \in W \qquad\quad \iff \qquad\quad \varphi_Q \in W,
\]
où $\pi \in W$ signifie que $Q$ est $W$-cocartésien, q.e.d.

\msp{page 265}

\dgrblock{dgrb598}{A four-object circle behind a strict pushout}
\begin{dcomment}{examples}{p264}{A four-object circle behind a strict pushout}
Return to the two-point span whose cylinder cofiber was calculated at ms.~230. The integral gives a smaller categorical model of the same circle, and explains why the word ``homotopy'' cannot be omitted from Proposition 20. Let $S$ be the discrete category on two objects $a,b$, and take $X_1=X_2=e$, both maps $S\to e$ being the unique functor. The ordinary pushout in $\mathrm{Cat}$ is $e$.

The fibred integral of this span has four objects, which we denote $u,v,a,b$. Its only nonidentity arrows are
\[
 u\longrightarrow a,\qquad u\longrightarrow b,
 \qquad v\longrightarrow a,\qquad v\longrightarrow b.
\]
There are no composable pairs of nonidentity arrows. Thus its nerve has four vertices, four nondegenerate edges, and no nondegenerate simplices of dimension at least two. Its realization is the four-edge cycle, homeomorphic to a circle. Equivalently, the cellular boundary map has rank three, so its first homology is $\mathbb Z$.

The comparison $\mathcal X_Q\to e$ therefore fails to be in $W_\infty$. The strict pushout square is not a homotopy pushout, and its actual homotopy pushout has the homotopy type of $S^1$. This four-edge circle and the earlier eight-edge cylinder cofiber both model the suspension of the two-point space. Neither original leg $S\to e$ is a Dwyer immersion: both identify the two objects. Replacing these legs before taking the ordinary pushout is precisely what preserves the missing circle. The example also fixes the variance of the integral without relying on a pictorial convention for its arrows.
\end{dcomment}


\section*{8\quad Appendice : Retour sur les axiomes W(7), W(7 bis) (pages 75, 76).}

(\footnote[135]{Reflexion provisoire, supplanté par XVI.}). On a vu (prop. 6, p. 78) que W(7 bis) implique W(7). Montrons une réciproque partielle. % [sic] Reflexion -> Réflexion, missing accent; supplanté -> supplantée, agreement (in footnote 135)

\dgrblock{dgrb599}{What the two converse gluing assertions require}
\begin{dcomment}{exposition}{p265:Montrons une réciproque partielle.}{What the two converse gluing assertions require}
Grothendieck wants to recognize a weak equivalence from the pieces of a decomposition. The familiar direction starts with equivalences on the common part and the two pieces and concludes that their homotopy pushouts are equivalent. His appendix asks for two converses: can the common part be recovered from the pieces and their union, and can the pieces be recovered from the common part and their union? These are substantive additions to homotopy-colimit invariance.

Write a square as $A\to B$, $A\to C$, with fourth vertex $P$, and write
\[
   C_e(u)=e\amalg^{\mathrm h}_{\operatorname{dom}(u)}
                    \operatorname{cod}(u)
\]
for the homotopy cofiber formed using the terminal object $e$. Whenever the square is a homotopy pushout, an equivalence $A\to B$ gives an equivalence $C\to P$. Grothendieck's conservative axiom asks for the converse as well. Applying that converse to the defining square of $C_e(u)$ says that a terminal homotopy cofiber detects an equivalence. This detection property is the extra ingredient in the proposition, through the full hypothesis W(7 bis).

It is useful to keep the three levels separate as the proof unfolds. Invariance of a homotopy pushout is formal once a coherent homotopy-colimit theory exists. The fact that an equivalence has terminal homotopy cofiber is equally formal. The converse, that terminal cofiber forces equivalence, is generally false for spaces; a complete example appears below. Thus the proposition under W(7 bis) and the later theorem for a general category of models must be read with different hypotheses. The last theorem drops an assumption that its third assertion still needs.
\end{dcomment}


\textbf{Proposition.} \textit{Supposons} W(7 bis) (axiome du carré cocartésien) \textit{satisfait, et plaçons nous sous les conditions de l'énoncé de} W(7), \textit{et plus généralement, dans le cas d'un diagramme cubique}
\[
\begin{tikzcd}[ampersand replacement=\&, row sep=small, column sep=small]
 \& X_2 \arrow[rr] \arrow[dd, "f_2" near end] \& \& X_3 \arrow[dd, "f_3"] \\
X_0 \arrow[ur] \arrow[rr, crossing over] \arrow[dd, "f_0"'] \& \& X_1 \arrow[ur] \& \\
 \& Y_2 \arrow[rr] \& \& Y_3 \\
Y_0 \arrow[ur] \arrow[rr] \& \& Y_1 \arrow[ur] \arrow[from=uu, crossing over, "f_1"' near end] \&
\end{tikzcd}
\]
\textit{où les deux carrés $X$ et $Y$ (face supérieure et inférieure du cube) sont $W$-cocartésiens. Considérons les conditions suivantes sur ce diagramme}
\begin{enumerate}[label={}]
\item[\textit{a)}] \textit{$f_0 \in W$.}
\item[\textit{b)}] \textit{$f_1$ et $f_2 \in W$.}
\item[\textit{c)}] \textit{$f_3 \in W$.}
\end{enumerate}
\textit{Alors on a les implications}
\[
a \text{ \textit{\&} } b \Longrightarrow c, \qquad\quad b \text{ \textit{\&} } c \Longrightarrow a.
\]
\textit{De plus, si $W$ est stable par facteurs directs (p. ex. $W$ fortement saturé), alors on a aussi}
\[
a \text{ \textit{\&} } c \Longrightarrow b,
\]
\textit{donc la conjonction} \texttt{[de]} \textit{deux quelconques parmi les conditions a) b) c) implique la troisième.}

\textsc{Démonstration.} Utilisant la factorisation
\printedpage{173}
% head: A. Grothendieck : Dérivateurs, Ch. XII, Édition des Universités de Montpellier II et Paris VII
\msp{page 266}
fonctorielle d'une flèche quelconque (ici les $f_i$) en $pi$, avec $i$ immersion ouverte de Dwyer et $p \in W$ (p. 229), on est remené au cas où les $f_i$ sont des immersions ouvertes de Dwyer. Mais pour toute flèche $g : S \longrightarrow T$ qui est dans $\Cof_W$, soient $Z(g)$ son mapping-cone % [sic] remené -> ramené, spelling
\[
\begin{tikzcd}[ampersand replacement=\&, row sep=small]
S \arrow[rr, "g"] \arrow[dr] \& \& T \arrow[dr, "\alpha(g)"] \& \\
 \& e \arrow[rr, "\bar{g}"] \& \& Z(g),
\end{tikzcd}
\]
et
\[
\alpha(g) : T = \mathrm{but}(g) \longrightarrow Z(g)
\]
la flèche canonique (\footnote[136]{C'est un peu vif, car il n'est pas clair que $X_1' \amalg_{X_0} X_2' \longrightarrow X_1 \amalg_{X_0} X_2$ soit \texttt{[dans]} $W$, sans utiliser déjà un des trois ingrédients de W(7) (savoir $f_0, f_1, f_2 \in W \implies f_3 \in W$). Il semble donc que notre argument ne marche que moyennant une hypothèse supplémentaire de cet ordre \dots Mais il y a le \textit{lemme de Brown,} cf. XVI \S3, notamment p. 20.}). Celle-ci est fonctorielle en $g$, de sorte que le diagramme cubique donné se prolonge en un double cube
\[
\begin{tikzcd}[ampersand replacement=\&, row sep=small, column sep=small]
 \& X_2 \arrow[rr] \arrow[dd, "f_2" near end] \& \& X_3 \arrow[dd, "f_3"] \\
X_0 \arrow[ur] \arrow[rr, crossing over] \arrow[dd, "f_0"'] \& \& X_1 \arrow[ur] \& \\
 \& Y_2 \arrow[rr] \arrow[dd, "\alpha_2" near end] \& \& Y_3 \arrow[dd, "\alpha_3"] \\
Y_0 \arrow[ur] \arrow[rr, crossing over] \arrow[dd, "\alpha_0"'] \& \& Y_1 \arrow[ur] \arrow[from=uu, crossing over, "f_1"' near end] \& \\
 \& Z_2 \arrow[rr] \& \& Z_3 \\
Z_0 \arrow[ur] \arrow[rr] \& \& Z_1 \arrow[ur] \arrow[from=uu, crossing over, "\alpha_1"' near end] \&
\end{tikzcd}
\]
avec
\[
Z_i = Z(f_i),\ \alpha_i = \alpha(f_i) \qquad\quad \text{pour } i \in \{0,1,2,3\}.
\]
Notons que W(7 bis) implique que si $g \in \Cof_W$,
\msp{page 267}
on a
\[
g \in W \qquad\quad \iff \qquad\quad \bar{g} \in W, \qquad\quad \text{i.e. } Z(g)\ W\text{-asphérique},
\]
en particulier
\[
f_i \in W \qquad\quad \iff \qquad\quad Z_i\ W\text{-asphérique}.
\]
Ainsi, l'énoncé à prouver sur le cube initial se ramène à un énoncé sur le carré
\printedpage{174}
% head: A. Grothendieck : Dérivateurs, Ch. XII, Édition des Universités de Montpellier II et Paris VII
\[
\begin{tikzcd}[ampersand replacement=\&, row sep=small, column sep=small]
 \& Z_1 \arrow[dr] \& \\
Z_0 \arrow[ur] \arrow[dr] \& \& Z_3, \\
 \& Z_2 \arrow[ur] \&
\end{tikzcd}
\]
savoir que les trois conditions
\begin{enumerate}[label={}]
\item[a$'$)] $Z_0$ $W$-asphérique.
\item[b$'$)] $Z_1$ et $Z_2$ $W$-asphériques.
\item[c$'$)] $Z_3$ $W$-asphérique.
\end{enumerate}
donnant lieu aux implications
\[
\mathrm{a}' \text{ \& } \mathrm{b}' \Longrightarrow \mathrm{c}', \qquad\quad \mathrm{b}' \text{ \& } \mathrm{c}' \Longrightarrow \mathrm{a}',
\]
et que si $W$ est stable par facteurs directs, alors
\[
\mathrm{a}' \text{ \& } \mathrm{c}' \Longrightarrow \mathrm{b}'.
\]

C'est à peu de choses près le corollaire de la page 76, complété à la page 77 (corollaire), à cela près qu'on y suppose les flèches des immersions ouvertes, et le carré cartésien (donc aussi cocartésien) dans Cat, i.e. $Z_3 = Z_1 \cup Z_2$, $Z_0 = Z_1 \cap Z_2$. Mais
\msp{page 268}
ce qui importe pour la validité de cet énoncé, c'est (on s'en doute) que $Z_*$ soit $W$-cocartésien. Donc il faut d'abord prouver ceci.

\dgrblock{dgrb600}{The gluing step that Brown supplies}
\begin{dcomment}{context}{p267:C'est à peu de choses près le corollaire de la page 76}{The gluing step that Brown supplies}
The factorization into a Dwyer immersion and a weak equivalence gives Grothendieck an explicit replacement of every vertical edge. To use the replacement on an entire cube, however, one must know that the induced replacement of a pushout is again a weak equivalence. Footnote~136 identifies this dependence. It is a gap in the chosen proof, not a counterexample to the conditional proposition.

The later discussion is unusually precise about the dependency. In Chapter~XVI, Commentaire~3.2, Grothendieck observes that the appendix had used the fact that Dwyer immersions are $W$-cofibrations at a point where his earlier proof of that fact already used a gluing assertion. Lemma~3.2.1 and Corollary~3.2.2 isolate the required Brown lemma; Proposition~3.3 then records the intended equivalence of the strong axioms. The chapter identifies the repair but also says explicitly that its author does not have a proof of that lemma there \citep[\S3, ms.~pp.~18--23]{GrothendieckDerivateursXVI}.

Here is the precise external result that closes this dependency. Suppose a class $C$ contains the isomorphisms and all maps from an initial object, is stable under composition and pushout, and every map factors as a member of $C$ followed by a weak equivalence. Suppose also that pushout preserves $C\cap W$. These are the relevant axioms of a category of cofibrant objects. Brown's gluing lemma implies that pushing any weak equivalence out along a member of $C$ preserves weak equivalences. In Brown's original formulation this is the dual of Part~I, \S4, Lemma~2; the proof uses his factorization lemma and the invariance result in Lemma~1 \citep[pp.~428--429]{Brown1973}. With $C$ the Dwyer immersions, this yields $C\subseteq\Cof_W$ and legitimizes the replacement of the cube. No converse detection property is supplied by Brown's lemma; that remains the role of W(7 bis).
\end{dcomment}


\textbf{Lemme.} \textit{Si le diagramme cubique p. 265 est tel que les faces $X$, $Y$ soient $W$-cocartésiennes, et les $f_i$ ($i \in \{0,1,2,3\}$) des $W$-cofibrations, alors le carré $Z$ des cofibres des $f_i$ est lui aussi $W$-cocartésien.}

Ce lemme résulte de la théorie du `déploiement des diagrammes cubiques', relatif à une catégorie de modèles $(\mathcal{M}, W)$ admettant une théorie des cofibres $W$-homotopiques, qu'il me faudra bien écrire! (Mais l'existence pour $(\mathrm{Cat}, W)$ d'une telle théorie des cofibres $W$-homotopiques, et par suite de celle des déploiements des diagrammes dans $\mathcal{M}$ de type $(\boldsymbol{\Delta}^1)^I$, $I$ ensemble d'indices fini, ne dépend que de W(7 bis) - on n'a pas à faire appel à W(7).)

\dgrblock{dgrb601}{Taking cofibers commutes with the pushout}
\begin{dcomment}{exposition}{p268:ne dépend que de W(7 bis)}{Taking cofibers commutes with the pushout}
The useful part of Grothendieck's large diagram is the compatibility of two successive homotopy pushouts. We can prove the lemma on p.~268 directly from this compatibility, without using the claim that a terminal cofiber detects an equivalence.

Let $J=(1\leftarrow0\to2)$, and let $X\to Y$ be a coherent morphism of homotopy-pushout squares in a homotopy theory with a terminal object and finite homotopy colimits. Put $Z_j=C_e(X_j\to Y_j)$. Computing the colimit in the two orders gives
\[
\begin{aligned}
 \mathop{\mathrm{hocolim}}_{j\in J} Z_j
 &\simeq
 \left(\mathop{\mathrm{hocolim}}_{j\in J}Y_j\right)
 \amalg^{\mathrm h}_{\mathop{\mathrm{hocolim}}_{j\in J}X_j}
 \left(\mathop{\mathrm{hocolim}}_{j\in J} e\right)\\
 &\simeq Y_3\amalg^{\mathrm h}_{X_3}e=Z_3.
\end{aligned}
\]
The first equivalence is the interchange of colimits of the single diagram indexed by $J$ and the cofiber span. The constant span $e\leftarrow e\to e$ has homotopy pushout $e$, which proves the second equivalence. Under these identifications the comparison map is the one in the given square $Z$, so this proves the lemma, including its compatibility with the prescribed edges. In derivator language the same calculation is the composition of left Kan extensions, with the parameter retained; the pointwise criterion is \citet[Corollary~3.14]{Groth2013}.

This proof also locates the excessive assertion on p.~269. From $u\in W$ one may conclude $C_e(u)\simeq e$ in any such theory. Going back from $C_e(u)\simeq e$ to $u\in W$ uses conservativity. The full W(7 bis) hypothesis permits that return, but the statement that the whole reduction used only its forward half is false. The pointwise-cofiber lemma survives without the reverse half; the reduction of all three cube implications does not.
\end{dcomment}


\msp{page 269}
Ce `deploiement' est un diagramme infini dont toutes les suites infinies, dans les trois directions des axes de coordonnées, sont des `suites de suspension'. Je reproduis la partie utile du déploiement. % [sic] deploiement -> déploiement, missing accent

% ---- from Der-ChXIIc-30-03-03_p061-064.tex
\printedpage{175}
% head: A. Grothendieck : Dérivateurs, Ch. XII, Édition des Universités de Montpellier II et Paris VII
% foot: Version 30 mars 2003

% dangling arrows (from X_2, X_3, Y_2, Y_3, Z_2, Z_3) as printed
\[
\begin{tikzcd}[ampersand replacement=\&, column sep=small, row sep=small]
 \& \& {} \& \& {} \& \& U \arrow[dd, "\mu"] \\
 \& X_2 \arrow[rr] \arrow[dd] \arrow[ur, shorten >=1ex] \& \& X_3 \arrow[rr] \arrow[dd] \arrow[ur, shorten >=1ex] \& \& T \arrow[ur] \arrow[dd] \& \\
X_0 \arrow[rr, crossing over] \arrow[ur] \arrow[dd] \& \& X_1 \arrow[rr, crossing over] \arrow[ur] \arrow[dd, crossing over] \& \& S \arrow[ur, "\lambda"] \arrow[dd, crossing over] \& \& U' \arrow[dd] \\
 \& Y_2 \arrow[rr] \arrow[dd] \arrow[ur, shorten >=1.5ex] \& \& Y_3 \arrow[rr] \arrow[dd] \arrow[ur, shorten >=1.5ex] \& \& T' \arrow[ur] \arrow[dd] \& \\
Y_0 \arrow[rr, crossing over] \arrow[ur] \arrow[dd] \& \& Y_1 \arrow[rr, crossing over] \arrow[ur] \arrow[dd, crossing over] \& \& S' \arrow[ur, "\lambda'"] \arrow[dd, crossing over] \& \& U'' \\
 \& Z_2 \arrow[rr] \arrow[ur, shorten >=1.5ex] \& \& Z_3 \arrow[rr] \arrow[ur, shorten >=1.5ex] \& \& T'' \arrow[ur] \& \\
Z_0 \arrow[rr] \arrow[ur] \& \& Z_1 \arrow[rr] \arrow[ur] \& \& S'' \arrow[ur, "\lambda''"] \& \&
\end{tikzcd}
\]

La condition que $X$ soit $W$-cocartésien s'exprime par $\lambda \in W$, ou encore par $U \in \Asph_W$ (car $U$ s'identifie à la cofibre homotopique $Z(\lambda)$ de $\lambda$ \dots\,), de même $Y$ $W$-cocartésien s'exprime par $U' \in \Asph_W$, et $Z$ $W$-cocartésien par $U'' \in \Asph_W$. Or $U'' \sim Z(\mu)$, où $\mu : U \longrightarrow U'$ est dans $W$ si $U$, $U'$ sont $W$-asphériques, donc $U'' = Z(\mu)$ est alors $W$-asphérique. Cela prouve le lemme.

Cela montre que, moyennant W(7 bis) (en fait, on a utilisé seulement la partie ‘universelle’ de W(7 bis)), savoir que dans
\begin{tikzcd}[ampersand replacement=\&, column sep=small, row sep=small, baseline=(\tikzcdmatrixname-2-1.base)]
T_0 \arrow[rr] \arrow[dr, "i"'] \& \& T_2 \arrow[dr, "i'"] \& \\
 \& T_1 \arrow[rr] \& \& T_3
\end{tikzcd}
diagramme cocartésien d'immersions ouvertes, $i \in W \Longrightarrow i' \in W$ (au lieu de $\Longleftrightarrow$),

\msp{page 270}

la validité de W(7) se ramène au cas où le carré $Y$ est le carré trivial
\begin{tikzcd}[ampersand replacement=\&, column sep=small, row sep=small, baseline=(\tikzcdmatrixname-2-1.base)]
 \& e \arrow[dr] \& \\
e \arrow[ur] \arrow[dr] \& \& e, \\
 \& e \arrow[ur] \&
\end{tikzcd}
$e$ l'objet final de Cat.

On peut maintenant conclure : On a (moyennant la partie directe de W(7 bis))
\[
\mathrm{a}' \ \& \ \mathrm{b}' \Longrightarrow \mathrm{c}',
\]
car si $Z_0$, $Z_1$ sont $\in \Asph_W$, alors $i_1 \in W$, donc $j_1 \in W$, et si de plus $Z_2 \in \Asph_W$, on a donc $Z_3 \in \Asph_W$.
\[
\begin{tikzcd}[ampersand replacement=\&, column sep=small, row sep=small]
 \& Z_1 \arrow[dr, "j_2"] \& \\
Z_0 \arrow[ur, "i_1"] \arrow[dr, "i_2"'] \& \& Z_3 \\
 \& Z_2 \arrow[ur, "j_1"'] \&
\end{tikzcd}
\]

Si on a aussi la partie inverse (‘inhabituelle’) de W(7 bis), alors
\[
\mathrm{b}' \ \& \ \mathrm{c}' \Longrightarrow \mathrm{a}',
\]
car si $Z_1$ et $Z_3$ sont $W$-asphériques, alors $j_2 \in W$, donc $i_2 \in W$, et si de plus $Z_2$ est $W$-asphérique, on conclut $Z_0$ $W$-asphérique.

\printedpage{176}
% head: A. Grothendieck : Dérivateurs, Ch. XII, Édition des Universités de Montpellier II et Paris VII
% foot: Version 30 mars 2003

Enfin, supposons que $W$ soit \textit{stable par facteurs directs,} et W(7 bis) satisfait sous forme directe. Je dis qu'alors
\[
\mathrm{c}' \ \& \ \mathrm{a}' \Longrightarrow \mathrm{b}',
\]
i.e.
\[
Z_0 \text{ et } Z_3 \ \in \Asph_W \qquad\quad \Longrightarrow \qquad\quad Z_1 \text{ et } Z_2 \ \in \Asph_W.
\]
Factorisant $i_1$ et $i_2$ en
\[
\begin{aligned}
Z_0 &\xrightarrow{i'_1} Z'_1 \xrightarrow{\alpha_1} Z_1, \\
Z_0 &\xrightarrow{i'_2} Z'_2 \xrightarrow{\alpha_2} Z_2,
\end{aligned}
\]
avec $i'_1$, $i'_2$ des immersions ouvertes de Dwyer, et $\alpha_1, \alpha_2 \in W$, on est ramené au cas où

\msp{page 271}

$i_1$, $i_2$ sont des \textit{immersions ouvertes de Dwyer,} et où le carré est cocartésien (ce qui implique qu'il est $W$-cocartésien). Faisant alors le cochangement de base par $Z_0 \xrightarrow{p_0} Z'_0 = e$, on trouve un homomorphisme de carrés cocartésiens
\[
\begin{tikzcd}[ampersand replacement=\&, column sep=small, row sep=small]
 \& Z_1 \arrow[rr] \arrow[dd, "p_1" near end] \& \& Z_3 \arrow[dd, "p_3"] \\
Z_0 \arrow[ur] \arrow[rr, crossing over] \arrow[dd, "p_0"'] \& \& Z_2 \arrow[ur] \& \\
 \& Z'_1 \arrow[rr] \& \& Z'_3, \\
e = Z'_0 \arrow[ur] \arrow[rr] \& \& Z'_2 \arrow[ur] \arrow[from=uu, crossing over, "p_2"' near end] \&
\end{tikzcd}
\]
où toutes les faces du cube sont $W$-cocartésiennes, et comme $p_0 \in W$ ($Z_0$ étant $W$-asphérique), les autres flèches verticales $p_1$, $p_2$, $p_3$ sont \texttt{[dans]} $W$, donc le carré $Z'_*$ satisfait aussi $Z'_0$ et $Z'_3$ \texttt{[}$W$\texttt{-]}asphériques (et de plus $Z'_0 = e$), et il faut prouver $Z'_1$, $Z'_2$ $W$-asphériques (d'où itou pour $Z_1$, $Z_2$). Donc, on est ramené au cas particulier où $Z_0 = e$.

Mais dans ce cas il existe une rétraction (triviale!) de $Z_1$ sur $Z_0$, donc aussi de $Z_3$ sur $Z_2$. Ainsi, $Z_2$ est un facteur

\msp{page 272}

direct d'un objet $W$-asphérique, ce qui implique, par l'hypothèse sur $W$, que $Z_2$ est $W$-asphérique. Par symétrie, $Z_1$ l'est aussi. On a gagné!

\dgrblock{dgrb602}{The retract argument is valid at the level of objects}
\begin{dcomment}{exposition}{p272:Par symétrie, $Z_1$ l'est aussi. On a gagné!}{The retract argument is valid at the level of objects}
Grothendieck's reduction to a common part equal to $e$ is sound and deserves to be retained. If
\[
 Z_3\simeq Z_1\amalg^{\mathrm h}_{Z_0}Z_2,
 \qquad Z_0\simeq e,
\]
then homotopy-colimit invariance identifies $Z_3$ with the homotopy coproduct of two objects under $e$. Each summand is a retract of that coproduct: the retraction onto $Z_1$ is induced by the identity of $Z_1$ and the composite $Z_2\to e\to Z_1$. Both agree on $e$, so the universal property supplies the retraction; its composite with the summand inclusion is the identity. The same construction applies to $Z_2$. Consequently, if $Z_3$ is $W$-aspherical and $W$-aspherical objects are closed under retracts, both summands are $W$-aspherical. This proves the object statement without any conservative axiom.

For the source's weaker hypothesis of closure under strict retracts, realize this argument using cofibration representatives $e\to Z_i^c$ in the category under $e$ and their ordinary pushout $P^c=Z_1^c\amalg_e Z_2^c$. These representatives compute the displayed homotopy pushout. The identity of $Z_1^c$ and the composite $Z_2^c\to e\to Z_1^c$ agree strictly on $e$, so the pushout retraction and its splitting are literal maps, with composite the identity. The comparison with $Z_3$ makes $P^c$ aspherical; strict retract closure applies to $Z_i^c$, and the replacement weak equivalences give the assertion for $Z_i$.

The phrase \emph{facteur direct} here means a retract, as the explicit construction shows. Strong saturation of $W$ also suffices because localization carries a retract diagram to a retract diagram, and a retract of an isomorphism is an isomorphism.

What this argument does not supply is the return from the objects $Z_i=C_e(f_i)$ to the maps $f_i$. Terminal cofibers need not detect equivalences. The object statement just proved therefore cannot by itself justify assertion~3\,\textdegree{} of the following general theorem, or the equivalence asserted in Corollary~1. Under the full conservative hypothesis that return is valid, and the final Corollary~2 survives.
\end{dcomment}


\textbf{NB} En fait, ces arguments sont généraux, et donnent ceci :

\setcounter{footnote}{136}
\textbf{Théorème.} \textit{Soit $(\mathcal{M}, W)$ une catégorie des modèles, admettant une théorie de la cofibre $W$-homotopique \textup{(}\footnote{Mais on ne sait vérifier l'existence d'une telle théorie dans $(\mathrm{Cat}, W)$ que si on admet la forme faible de W(7), i.e. $f_0, f_1, f_2 \in W \implies f_3 \in W$, ou quelque cas particulier de celle-ci. Mais le \textit{lemme de Brown} sauve la situation, cf. XV \S3.}\textup{)} (i.e. on a une factorisation de toute flèche $f : X \longrightarrow Y$ en $pi$, avec $i \in \Cof_W$, $p \in W$.) On a donc une notion de carré $W$-cocartésien dans $\mathcal{M}$ \textup{(}\footnote{à vérifier que la factorisation suffit, même si tous les objets ne sont pas fibrants sur $e$. Mais on peut alors remplacer $\mathcal{M}$ par $\mathcal{M}_f$, et développer ce formalisme des fibres homotopiques en conséquence \dots}\textup{)}. Soit
\printedpage{177}
% head: A. Grothendieck : Dérivateurs, Ch. XII, Édition des Universités de Montpellier II et Paris VII
% foot: Version 30 mars 2003
alors $f : X \longrightarrow Y$ un homomorphisme d'un carré $W$-cocartésien $X$ dans un autre $Y$, $f = (f_0, f_1, f_2, f_3)$.}
\[
\begin{tikzcd}[ampersand replacement=\&, column sep=small, row sep=small]
 \& X_2 \arrow[rr] \arrow[dd, "f_2" near end] \& \& X_3 \arrow[dd, "f_3"] \\
X_0 \arrow[ur] \arrow[rr, crossing over] \arrow[dd, "f_0"'] \& \& X_1 \arrow[ur] \& \\
 \& Y_2 \arrow[rr] \& \& Y_3 \\
Y_0 \arrow[ur] \arrow[rr] \& \& Y_1 \arrow[ur] \arrow[from=uu, crossing over, "f_1"' near end] \&
\end{tikzcd}
\]

\textit{Considérons les conditions sur $f$}

\begin{enumerate}
\item[\textit{a)}] \textit{$f_0 \in W$.}
\item[\textit{b)}] \textit{$f_1$ et $f_2 \in W$.}
\item[\textit{c)}] \textit{$f_3 \in W$.}
\end{enumerate}

\textit{Alors}

\begin{enumerate}
\item[\textit{1°)}] \textit{a \& b $\Longrightarrow$ c.}

\msp{page 273}

\item[\textit{2°)}] \textit{Supposons que $\mathcal{M}$ ait un objet final $e$. Alors l'implication}
\[
\textit{b \& c} \Longrightarrow \textit{a}
\]
\textit{pour}~tout \textit{homomorphisme $f : X \longrightarrow Y$ de carrés $W$-cocartésiens équivaut à l'axiome suivant sur $W$ :}

Axiome des carrés $W$-cocartésiens conservatifs (\footnote{cf. argument page 78 qui montre que l'implication b \& c $\Longrightarrow$ a entraine l'axiome conservatif.}) : \textit{Si}
\[
\begin{tikzcd}[ampersand replacement=\&, column sep=small, row sep=small]
 \& Z_1 \arrow[dr, "j_2"] \& \\
Z_0 \arrow[ur, "i_1"] \arrow[dr, "i_2"'] \& \& Z_3 \\
 \& Z_2 \arrow[ur, "j_1"'] \&
\end{tikzcd}
\]
\textit{est $W$-cocartésien, alors $j_1 \in W \Longrightarrow i_1 \in W$ (donc en fait $i_1 \in W \Longleftrightarrow j_1 \in W$ et $i_2 \in W \Longleftrightarrow j_2 \in W$).}

\textit{Cet axiome est équivalent à la conjonction des deux axiomes suivants.}

\dgrblock{dgrb603}{A direct test for conservative homotopy pushouts}
\begin{dcomment}{exposition}{p273:Cet axiome est équivalent}{A direct test for conservative homotopy pushouts}
Grothendieck now isolates the property needed by the preceding reduction. Once homotopy pushouts are available coherently, there is a short proof of the test that also clarifies Cons~1 and Cons~2. We work in a homotopy theory with a terminal object and finite homotopy colimits, and use ``equivalence'' for a morphism detected as an isomorphism by its localization.

For every homotopy-pushout square
\[
\begin{tikzcd}[ampersand replacement=\&]
 A\arrow[r,"u"]\arrow[d] \& B\arrow[d]\\
 C\arrow[r,"v"'] \& P
\end{tikzcd}
\]
associativity gives a canonical equivalence
\[
 C_e(v)=e\amalg_C^{\mathrm h}P
 \simeq e\amalg_C^{\mathrm h}(C\amalg_A^{\mathrm h}B)
 \simeq e\amalg_A^{\mathrm h}B=C_e(u).
\]
Thus, if terminal homotopy cofiber detects equivalences, $v$ being an equivalence forces $u$ to be one. Conversely, apply the conservative-pushout axiom to the defining cofiber square: if $C_e(u)\simeq e$, its lower edge $e\to C_e(u)$ is an equivalence, and the axiom forces $u$ to be an equivalence. This proves the equivalence of conservativity and Cons~1, with the latter interpreted for a cofibration replacement of every map. The intended cocartesian Cons~2 follows: if $B,C,P$ are terminal up to equivalence, the edge $B\to P$ is an equivalence, hence so is $A\to C$; consequently $A\simeq e$. It is therefore redundant under this coherent formulation.

These observations complete the proof of the two converse cube implications under the appropriate hypothesis. Given a morphism $f:X\to Y$ of homotopy-pushout squares, its pointwise cofibers $Z_i=C_e(f_i)$ form a homotopy-pushout square by the preceding lemma. If $f_1,f_2,f_3$ are equivalences, then $Z_1,Z_2,Z_3\simeq e$; conservativity gives $Z_0\simeq e$, and Cons~1 gives $f_0\in W$. Conversely, the implication $b\mathbin{\&}c\Rightarrow a$ applied to the map from a square $(A,B,C,P)$ to $(B,B,P,P)$ with components $(u,1_B,v,1_P)$ proves conservativity immediately.

If instead $f_0,f_3$ are equivalences and aspherical objects are closed under retracts, the object argument on pp.~271--272 gives $Z_1,Z_2\simeq e$; Cons~1 then gives $f_1,f_2\in W$. Together with ordinary pushout invariance this proves the final Corollary~2. It also identifies exactly what must be added to the printed assertion~3\,\textdegree{}: a cofiber-detection hypothesis, equivalently conservative homotopy pushouts, in addition to the retract condition.
\end{dcomment}


\textbf{Cons 1.} \textit{Si on a un carré cocartésien}
\printedpage{178}
% head: A. Grothendieck : Dérivateurs, Ch. XII, Édition des Universités de Montpellier II et Paris VII
% foot: Version 30 mars 2003
\[
\begin{tikzcd}[ampersand replacement=\&, column sep=small, row sep=small]
S \arrow[rr] \arrow[dr] \& \& T \arrow[dr] \& \\
 \& e \arrow[rr] \& \& U,
\end{tikzcd}
\]
\textit{avec $f \in \Cof_W$, alors $U \in \Asph_W \Longrightarrow f \in W$ (donc $U \in \Asph_W \Longleftrightarrow f \in W$).}

\textbf{Cons 2.} \textit{Si on a un carré cartésien,}
\[
\begin{tikzcd}[ampersand replacement=\&, column sep=small, row sep=small]
 \& Z_1 \arrow[dr] \& \\
Z_0 \arrow[ur, "i_1"] \arrow[dr, "i_2"'] \& \& Z_3, \\
 \& Z_2 \arrow[ur] \&
\end{tikzcd}
\]
\textit{avec $i_1, i_2 \in \Cof_W$, alors : Si $Z_1$, $Z_2$, $Z_3$ sont $W$-asphériques, $Z_0$ l'est aussi.}

\dgrblock{dgrb604}{Why Cons 2 must concern pushouts}
\begin{dcomment}{examples}{p273:Si on a un carré cartésien}{Why Cons 2 must concern pushouts}
The edition prints \emph{cartésien} in Cons~2. Taken literally, this gives a different and false consequence of conservativity. The issue can be seen inside a stable homotopy theory, where conservative homotopy pushouts do hold.

Let $k$ be a field and use homologically graded chain complexes. Put $A=ka$ in degree zero, and let
\[
 B=(kb\xrightarrow{d}ka),\qquad
 C=(kc\xrightarrow{d}ka),\qquad d(b)=d(c)=a.
\]
Embed both in the complex
\[
 D=(kd\longrightarrow kb\oplus kc\longrightarrow ka),
 \qquad \partial d=b-c,\quad \partial b=\partial c=a.
\]
Here $d$ as a basis symbol lies in degree two. Both $B$ and $C$ are acyclic. In $D$ the first differential is injective, its image is the kernel of the second, and the second is surjective, so $D$ is acyclic too. Their ordinary pullback over $D$ is the intersection $A$, whose zeroth homology is $k$. All four inclusions are degreewise monomorphisms. Such a map is a $W$-cofibration for quasi-isomorphisms: pushout leaves its cokernel unchanged, and the long exact homology sequences show that a weak equivalence of sources induces one of the pushouts. Over a field these maps are also cofibrations in the usual model structure.

Thus a strict cartesian square with the printed cofibration hypotheses need not satisfy Cons~2, even when the surrounding derived theory is stable. Replacing \emph{cartésien} by \emph{cocartésien} restores the statement used in the argument. The source wording is retained in the text and recorded as an edition slip.
\end{dcomment}


\msp{page 274}

\item[\textit{3°)}] \textit{Supposons que tout facteur direct d'un objet $W$-asphérique soit $W$-asphérique, ce qui est le cas notamment si $W$ est lui-même stable par facteurs directs, a fortiori si $W$ est fortement saturé. Supposons de plus que $\mathcal{M}$ ait} \texttt{[un]} \textit{objet final. Alors on a (sous les conditions générales du début)}
\[
\textit{c \& a} \Longrightarrow \textit{b,}
\]
\textit{i.e. si $f_0$, $f_3$ sont} \texttt{[dans]} \textit{$W$, de même $f_1$, $f_2$.}
\end{enumerate}

Cette dernière partie peut se déployer ainsi :

\dgrblock{dgrb605}{An explicit cube testing the generalization}
\begin{dcomment}{development}{p274:Cette dernière partie peut se déployer ainsi}{An explicit cube testing the generalization}
Grothendieck's retraction proof invites a test of the stronger assertion about arbitrary maps of squares. We give the test in full, since it distinguishes a missing hypothesis from a missing step in a proof. The acyclic complex used below is classical; the assembly into the two diagnostic cubes is our application to this passage. We checked the source's later Chapter~XVI discussion, Brown's gluing lemmas, and the cofiber and Mayer--Vietoris treatments cited here; we did not find this particular counterexample to the printed generalization in those sources. No priority is claimed.

\textbf{A cube with $a$ and $c$, but not $b$.}\ours\ 
In simplicial sets with weak homotopy equivalences, there is a map between homotopy-pushout squares whose zeroth, second and third components are equivalences and whose first component is not.

\emph{Proof.} Begin with the finite CW complex $A$ having one vertex, two edges $a,b$, and two two-cells attached by $a^5b^{-3}$ and $b^3(ab)^{-2}$ \citep[Example~2.38]{Hatcher2002}. Its cellular boundary in positive dimensions is
\[
 0\longrightarrow\mathbb Z^2
 \xrightarrow{\left(\begin{smallmatrix}5&-2\\-3&1\end{smallmatrix}\right)}
 \mathbb Z^2\longrightarrow0.
\]
The determinant is $-1$, so $A$ is acyclic. Its fundamental group is nontrivial: sending $a$ to $(12345)$ and $b$ to $(153)$ respects the relations, since their product is $(23)(45)$, and gives a nontrivial permutation representation. The suspension $\Sigma A$ is simply connected by van Kampen and acyclic by the suspension isomorphism in reduced homology. Hurewicz and Whitehead therefore make it contractible \citep[Theorem~4.32 and Corollary~4.33]{Hatcher2002}.

Choose a simplicial set $K$ of the homotopy type of $A$, let $CK=K\star\Delta^0$, and write $e=\Delta^0$. The inclusion $K\hookrightarrow CK$ is a cofibration, $CK$ is contractible, and $CK\amalg_K e$ has the homotopy type of $\Sigma A$, hence is weakly contractible. Consider
\[
 X=\left(\begin{tikzcd}[ampersand replacement=\&, column sep=small,row sep=small]
 K\arrow[r,equal]\arrow[d,hook] \& K\arrow[d,hook]\\
 CK\arrow[r,equal] \& CK
 \end{tikzcd}\right),\qquad
 Y=\left(\begin{tikzcd}[ampersand replacement=\&, column sep=small,row sep=small]
 K\arrow[r]\arrow[d,hook] \& e\arrow[d]\\
 CK\arrow[r] \& CK\amalg_K e
 \end{tikzcd}\right).
\]
Both strict pushouts compute homotopy pushouts, since the common vertical leg is a cofibration. There is a strict map $X\to Y$ with components
\[
 (f_0,f_1,f_2,f_3)=(1_K,\ K\to e,\ 1_{CK},\ CK\to CK\amalg_K e).
\]
Its zeroth and second components are identities; its third connects weakly contractible objects; its first is not a weak equivalence because $\pi_1(K)\ne1$. This proves the assertion. All simplicial sets are cofibrant, factorizations exist, and weak equivalences are strongly saturated and closed under retracts. Thus the failure is present under the printed hypotheses of assertion~3\,\textdegree{}, not a failure of replacement machinery.

The same $K$ tests the other converse. Map the homotopy-pushout square with vertices $(K,CK,CK,CK\amalg_K CK)$ to the constant square at $e$. Its two cone vertices and its suspension vertex are weakly contractible, while $K$ is not. Hence $b\mathbin{\&}c\Rightarrow a$ fails for spaces, as the conservative axiom predicts. Finally, the object statement in Corollary~1 remains valid for spaces by the retract proof above. Its proposed equivalence with $a\mathbin{\&}c\Rightarrow b$ therefore fails as printed. Adding cofiber detection repairs both assertions.
\end{dcomment}


\textbf{Corollaire 1.} \textit{Sous l'hypothèse supplémentaire} \textup{3\,°)} \textit{sur $W$, la validité de c \& a $\Longrightarrow$ b (pour toute flèche $f$ entre carrés $W$-cocartésiens) équivaut à chacun des deux autres axiomes suivants sur $W$ (qui sont donc équivalents) :}

\begin{enumerate}
\item[\textit{1)}] \textit{Si on a un carré $W$-cocartésien}
\[
\begin{tikzcd}[ampersand replacement=\&, column sep=small, row sep=small]
 \& Z_1 \arrow[dr] \& \\
Z_0 \arrow[ur] \arrow[dr] \& \& Z_3, \\
 \& Z_2 \arrow[ur] \&
\end{tikzcd}
\]
\textit{alors $Z_0$ et $Z_3$ $W$-asphériques $\Longrightarrow$ $Z_1$ et $Z_2$ $W$-asphériques.}
\item[\textit{1$'$)}] \textit{Comme 1), avec $Z_0 = e$ (dans le cas où on suppose que $\mathcal{M}$ admet} \texttt{[un]} \textit{objet final $e$).}
\end{enumerate}

\msp{page 275}

\textbf{Corollaire 2.} \textit{Supposons que $W$ satisfasse l'axiome des carrés $W$-cocartésiens conservatifs, que tout facteur direct d'un objet $W$-asphérique soit $W$-asphérique (p. ex. $W$ stable par facteurs directs), enfin que $\mathcal{M}$ admette un objet final. Alors pout tout homomorphisme $f : X \longrightarrow Y$ de carrés $W$-cocartésiens, la conjonction de deux parmi les conditions a) b) c) de p. 272 entraîne la troisième.} % [sic] pout -> pour, typographical error


\dgrblock{dgrb606}{The stable form of the cube calculus}
\begin{dcomment}{later}{p275}{The stable form of the cube calculus}
The last corollary has a particularly transparent home in a stable derivator. Coherent cocartesian squares then are also cartesian, and their geometry gives a Mayer--Vietoris triangle
\[
 X_0\xrightarrow{(i_1,-i_2)}X_1\oplus X_2
 \xrightarrow{(j_2,j_1)}X_3\longrightarrow\Sigma X_0.
\]
The existence of this triangle is \citet[Theorem~6.1, arXiv v2]{GrothPontoShulman2014}; their Theorem~7.1 compares stability with the invertibility of suspension and of the cofiber/fiber adjunction. This gives an established realization of the suspension calculus suggested by the large diagram on p.~269.

For a morphism of two such squares, the induced map of triangles makes the three implications immediate but still instructive. Equivalences on the first and middle terms give one on the third; equivalences on the middle and third give one on the first. Equivalences on the first and third give an equivalence
$f_1\oplus f_2$. Each $f_i$ is a retract of that direct-sum map in the arrow category, so each is an equivalence. This is the additive form of the retraction argument, with cofiber detection now supplied by stability.

Conservativity should nevertheless not be identified with stability. The former is the particular detection property isolated on p.~273; stability also requires a pointed theory and the agreement of cartesian and cocartesian squares. The source's general constructions of coherent squares and iterated cofibers remain useful in unstable theories, while the converse arguments require their stated additional hypotheses.
\end{dcomment}
\begin{otherlanguage}{english}
\dgroverview{XIIc-b}{Chapter XII \textperiodcentered\ \textup{XIIc-b} \textperiodcentered\ Overview}
% Overview by GPT-6 Astra (OpenAI, ChatGPT Pro mode).
\begingroup\selectlanguage{english}
\noindent\textbf{From replacements to a calculus of coherent squares.}
The pages gathered here form a single mathematical movement. Grothendieck first constructs categorical cylinders and uses them to establish the gluing rules for homotopy pushouts. He then asks which part of this geometry can be recovered from a category $\mathcal M$ and a class $W$ of weak equivalences alone. The answer takes the form of a homotopy completion of a cospan, characterized by an adjunction. Finally he returns to pushouts and investigates how far weak equivalences can be recovered from the pieces of a decomposition. The first two movements build the calculus; the appendix tests a considerably stronger conservativity property.

These pages are manuscript pages 228--275, printed pages 145--178 of the edition \citep{Grothendieck1990XIIc}: the second part of section~VI of Chapter~XII, that is, the end of its section~6 followed by sections~7 and~8. The audit covers that entire range, including the fragment of a footnote continued at the beginning. Page labels in the margins reproduce the two numbering systems of the transcription; citations in the commentary use manuscript pages unless another system is specified.

\medskip\noindent\textbf{Why the diagrams must precede localization.}
The motivation is older than the technical apparatus of these pages. In section 69 of \emph{Pursuing Stacks}, Grothendieck traces the need for derivators to the trace formalism of SGA~5: a single derived category did not retain the diagrammatic information needed for the desired compatibilities \citep[\S69, pp.~1--2]{Grothendieck1983}. His letter to Thomason of 2 April 1991 returns to this history and to the system of localized diagram categories \citep[\S2, pp.~5--6]{Grothendieck1991Letter}. These are historical statements checked in those texts; we do not infer that every later definition was already present there in its final form.

Here this motivation becomes a concrete distinction. With the manuscript's contravariant convention, the relevant category is
\[
 \mathbb D(I)=\operatorname{Fun}(I^{\mathrm{op}},\mathcal M)
                  [W(I)^{-1}],
\]
where $W(I)$ consists of the objectwise weak equivalences. It is not obtained by first passing to $\mathcal M[W^{-1}]$ and then forming its ordinary diagram category. A cospan remembers two maps to a common base together; its homotopy completion must preserve that relationship and its naturality in all three vertices. This explains both the attention given to the whole completed square and the two localization lemmas whose proofs the manuscript postpones. The modern introduction in \citet[pp.~313--315]{Groth2013} makes the same structural distinction.

\medskip\noindent\textbf{The constructive part and its later justification.}
The fence interval $\Phi_n$, with $n\geq2$, lets both endpoints enter by a Dwyer immersion. The associated mapping cylinder gives a functorial factorization into a $W$-cofibration followed by a weak equivalence. The commentary proves the endpoint condition and follows the construction on the map from two discrete objects to a point. Taking its categorical cofiber produces an eight-edge circle. Later, the Grothendieck integral of the same span gives a four-edge circle. These are two explicit models of the homotopy pushout; the ordinary categorical pushout is only a point.

The appeal to Thomason is valid in the needed cellular form. It must not be strengthened to the assertion that every Thomason cofibration is a Dwyer map: the retract assertion used to obtain that conclusion is false, and \citet{Cisinski1999Dwyer} gives the counterexample and repair. The recent comparison of Dwyer pushouts with $(\infty,1)$-categorical pushouts is stronger in a different direction, identifying what the nerve preserves rather than adding the appendix's conservative axiom \citep[Theorems~1.5--1.6, arXiv v3]{HackneyOzornovaRiehlRovelli2024}.

In section 7, a $W$-fibration means a map whose pullback functor preserves the designated weak equivalences. The factorization assumption says that every arrow is a weak equivalence followed by such a fibration. It is essential to retain this intrinsic meaning: the manuscript has not assumed a Quillen model structure with arbitrary lifting classes. Under its stated terminal-object and pullback hypotheses, these data form a left derivable category. The commentary verifies the nontrivial axiom, stability of trivial fibrations under arbitrary base change, by a two-step factorization and two-out-of-three argument.

This observation gives a complete modern repair of the deferred Lemmas 1 and 2. We construct the appropriate fibration structures on the good cospans and good squares, check the factorization and pullback conditions, and verify approximation for their inclusions in the ambient diagram categories. Cisinski's Proposition~1.8 also supplies the strengthened replacement assertion for every full intermediate subcategory of a slice containing its fibrant objects; his Theorem~3.12 establishes the two diagram-localization equivalences \citep{Cisinski2010}. No functorial factorization or global assumption that all objects are fibrant is required. A separate, shorter deformation argument explains what functorial factorization would give. The audit still labels the source proofs as deferred or sketched: the later completion does not change what is written on these pages.

With those equivalences available, restriction from squares to cospans has a fully faithful right adjoint. On a cospan $A\to S\leftarrow B$ with one good leg, its missing corner is represented by $A\times_S B$. The groupoid example $*\to BG\leftarrow *$ gives the discrete groupoid $G$ at that corner, whereas the strict pullback without replacement is a point. The adjunction retains the two projections and their compatibility, as well as the corner's homotopy type. Proposition 19 then gives an intrinsic test for cartesianness, and we prove its independence of choices, invariance under objectwise weak equivalences, and pasting directly from the factorization axiom. Strong saturation enters precisely when an isomorphism after localization must be recognized as an original member of $W$.

\medskip\noindent\textbf{What the appendix adds, and where its generalization fails.}
The ordinary gluing implication says that equivalences on the common part and the two pieces induce an equivalence on their homotopy pushouts. The appendix asks for converses. Its full W(7 bis) hypothesis is substantially stronger: a homotopy pushout must also reflect weak equivalences in an attached leg. In a coherent homotopy-pushout theory this is equivalent to detection by the terminal homotopy cofiber
\[
 C_e(u)=e\amalg^{\mathrm h}_{\operatorname{dom}(u)}
                         \operatorname{cod}(u).
\]
The forward implication $u\in W\Rightarrow C_e(u)\simeq e$ is formal. The reverse implication is the additional detection property. We prove the equivalence with conservativity by pasting cofiber squares.

This distinction exposes an actual error in the general theorem on manuscript page 274. Retract closure of aspherical objects alone does not imply that equivalences at the initial and final vertices of a map of homotopy-pushout squares force equivalences at the two middle vertices. Our counterexample uses Hatcher's finite acyclic noncontractible complex $A$. Its suspension is contractible; the map from the square $(A,A,CA,CA)$ to $(A,*,CA,CA/A)$ has three weak-equivalence components and one component $A\to *$ that is not a weak equivalence. Both squares are homotopy pushouts. We give the presentation of $A$, its cellular homology calculation, a nontrivial permutation representation of its fundamental group, and the suspension argument \citep[Example~2.38]{Hatcher2002}. The related equivalence asserted in Corollary 1 therefore fails as printed. Adding cofiber detection repairs both statements, and the final Corollary 2 already has the required conservative hypothesis.

There is also a source slip in Cons~2: the edition prints \emph{cart\'esien} where its argument requires \emph{cocart\'esien}. A finite chain-complex example shows that the literal cartesian assertion fails even in a stable theory. The proof of the intended criterion shows that its cocartesian version follows from Cons~1 under the coherent formulation. In stable derivators the entire corrected cube calculus has a familiar expression through the Mayer--Vietoris triangle \citep[Theorems~6.1 and 7.1, arXiv v2]{GrothPontoShulman2014}.

The appendix's earlier circularity is a different issue. Chapter XVI identifies the dependence on the Dwyer-cofibration gluing lemma but explicitly does not give its proof there \citep[\S3, ms.~pp.~18--23]{GrothendieckDerivateursXVI}. Brown's published gluing lemma, dualized, supplies the missing result \citep[Part~I, \S4, Lemma~2]{Brown1973}. Thus the initial conditional proposition survives, while the later hypothesis-free generalization does not. These are distinct diagnoses, with distinct repairs.

\medskip\noindent\textbf{Editorial and verification limits.}
The text preserves slips printed in the edition. Its one mathematical transcription correction restores the interior natural transformation $\alpha$ on manuscript page 235, which the transcription had drawn as an arrow between the wrong objects. The separate registers record the source's proof status, the present repair, and the evidence for each reading. Bibliography notes identify exactly which statements, proofs, or metadata were checked; in particular, the original proof of Thomason's 1979 theorem was not inspected, although its statement and a later proof were checked.

\endgroup
\end{otherlanguage}
\dgrchunkend
